% Modernised reading — fonds Alexandre Grothendieck, shelfmark 84, pages 1-192.
%
% This is an INTERPRETATION, not a transcription. It re-reads the pages in
% today's notation and vocabulary, states the hypotheses the manuscript leaves
% implicit, and moves everything the critical apparatus carried into
% footnotes. It is held to being mathematically correct as it stands; where
% the manuscript is loose or mistaken the true statement appears and the
% footnote says what the page has. In any dispute of substance the
% transcription governs: whoever cites Grothendieck cites batch-01.fr.tex to
% batch-10.fr.tex.
%
% Pass: Opus 5.5 (claude-opus-5-5), 2026-10-10 — first pass, unchecked against the pages by a human.
%
% Source: the ten batches as revised on the 700 dpi tiles on 2026-10-10 (each
% batch header carries its « % Revised: » line). Read whole, twice, before any
% of this was written.
%
% Scope. 192 facsimile pages, 152 transcribed. Absent from the transcription,
% and from this reading: pp. 4, 6, 7 (unrelated typescript and a computer
% listing), p. 44 (blank), p. 122 (blank), the even pages 124-140, 142-160 and
% 182-192 (typescript pages of his own non-mathematical writing used as scrap
% paper — the batch notes identify those of 142-160 as pages of Récoltes et
% Semailles), and the even pages 162-180 (blank). Page 2 carries a short
% personal note, not transcribed and not mentioned further here. Page 42 is a
% cover sheet in his hand, on music paper, numbered « (79) »; page 25 carries
% his own number « 30 ». Nothing of the typescripts is restated (RIGHTS.md).
%
% The group. Filed in « Géométrie et topologie combinatoire » (69-102); the
% subject is algebra over a general base. It meets folder 82 (SO(3) ≃ GP(1),
% dated [vers 1982]) on three points — the commutant of a non-scalar 2x2
% endomorphism and its characteristic-2 automorphism u -> u + Tr u (82,
% pp. 36-40; here pp. 14-15, 34, 38), the rank-4 Azumaya algebra with its
% reduced norm and canonical involution (82, pp. 46-55; here pp. 27-33) — and
% folder 98 on the norm of an invertible module along a finite locally free
% morphism (here p. 24). Cross-references are made at those places only.
%
% Spine. Three parts and an opening. 0 (pp. 1-5): leaves that presuppose
% Part II (a Z/2-graded monoid built from the functors of pp. 53-54; a
% triquadratic algebra; norm-one elements). I (pp. 8-41) « Formes
% quadratiques binaires et revêtements quadratiques »: similarity structures,
% the ring of homotheties of a rank-2 module, the algebra A_Q, the
% « Théorème récapitulatif » (binary forms ↔ quadratic algebra + invertible
% module), the Azumaya reformulation, symmetries. II (pp. 42-121, under the
% cover « χ-scindages, produits contractés … »): quadratic covers described
% by (L, P_1, Q) and, when 2 is regular, by (L, δ, T_0); characteristic 2;
% the contracted product X_1 * X_2 by way of χ-trivialised torsors; the
% coinvariant criterion; δ = δ_1 δ_2; relative trace and norm of
% A_1 ⊗ A_2 over A; quasi-quadratic and quasi-biquadratic algebras. III
% (pp. 123-191): the χ-formalism for extensions 0 -> L -> E -> M -> 0,
% pushout and pullback products, the symmetric operator F_11, and the pairing
% A × E -> O on a product of quadratic algebras.
%
% Notation fixed for the whole folder. X a scheme or locally ringed topos,
% O its structure sheaf. Quad(E) = Hom(Γ^2 E, O). A quadratic algebra A
% (commutative, locally free of rank 2), Tr, N, x̄ = σ(x) = Tr(x) - x;
% X' = Spec A its quadratic cover. Part II: E = A^∨, 0 -> L -> E -ψ-> O -> 0,
% ψ(x) = x(1), P_λ = ψ^{-1}(λ), T ∈ Γ P_2 the trace, N ∈ Γ Sym^2 E the norm;
% in a splitting x_0 ∈ Γ P_1, Q(x_0 + u) = u^2 + b u + c with
% b = π(x_0) = 2x_0 - T. Discriminants: δ = b^2 - 4c = T^2 - 4N (his δ,
% « le discriminant changé de signe », p. 43); d(Q) = det of the polar form
% (his δ(Q) of p. 18), so δ = -d. The product of χ-trivialised extensions
% (his ∧ p. 68, *, ⊛ p. 143, ∨ in Part III) is E_1 * E_2 throughout; the
% pullback product of Part III (his ∧, pp. 167-169) is E_1 *^ϖ E_2. The
% Azumaya algebra of § 4 is \mathcal{M} (his M), its anti-commutant
% {θ : θu = ūθ} is \mathcal{J} (his M of p. 39).
%
% Substantive departures from the pages, each footnoted where it occurs:
%  - p. 1: the « R_* = R ×_Z Z/2Z » of the box is read as the monoid of pairs
%    (n, x), x ∈ R_{n mod 2}; p. 2's (LL', 2δδ') read (LL', δδ'), per p. 1's
%    definition; the exponents of the Arf/c_1 formulas are uncertain and the
%    remark « (si 2 inversible) » cannot be meant as written.
%  - p. 3: « e_i inversibles i.e. A_i étales » needs 2e_i invertible.
%  - p. 8: Quad(E) = Hom(Γ^2 E, O) (page: Hom(Γ^2 E, E)); pp. 9-10: Sym^2(E^∨)
%    (page: Sym^2(E)).
%  - p. 13: with the conventions stated, O -> End(E) is 1 -> -1_E; as
%    submodules nothing changes.
%  - p. 16: Δ' ≅ (A/O)^∨ ⊗ L (page: A/O ⊗ L^∨, which is Δ).
%  - p. 17: Q_u(x) = x ∧ u(x) (box: u(x) ∧ x; his coordinates follow x ∧ u(x)).
%  - p. 18: d(Q) ∈ Γ((Δ' ⊗ L^∨)^{⊗2}) (an earlier line: L^{⊗2}).
%  - p. 21: « structure fidèle » in 3° d) read strictly faithful.
%  - p. 26: étale topology suffices for U_A-torsors only when U_A is smooth.
%  - p. 27: Azumaya algebra of rank 4 over O (page: « sur E »).
%  - pp. 34-38: the uniqueness of the conjugation holds up to the identity
%    near points where 2 = 0; stated precisely. Pages 35-38 are out of order;
%    the reading follows the order of the argument (34, 38, 37, 36, 35, 39).
%  - p. 37 « Norme 1 » vs pp. 39-40 det θ = -1: the reading follows 39-40.
%  - p. 40 c): lines of E (page: of A).
%  - p. 46 vs p. 52: Q(φ_N(x)) = -N(x)δ (p. 52 writes +; p. 97 has -).
%  - p. 63: 4c_1c_1 read 4c_1c_2. p. 75: α(x_1, x_2) read f(x_1, x_2).
%    p. 79 (margin): T = 2χx - b read χx_0 - b.
%  - p. 91: « X_1 × X_2 -> X_1 ∧ X_2 fid. plat » is false without a
%    hypothesis; a counter-example is given.
%  - p. 93: the minimal equation is u^2 + <u,b> u + c(u) = 0.
%  - pp. 111-112: T_1 read T_2, t_1(x_1) read n_1(x_1), k -> A_1 read k -> A_2.
%    p. 114: x_i = σ_i(x) (page: T_i(x)). p. 119, Lemma b): the ring of 𝒳,
%    not of X_J.
%  - p. 131: b_2'' = b_2' + χu_2'; the signs b_i' = b_i ± χu_i fixed as +.
%  - p. 165: F_12 = ψ_1 ⊗ ψ_2 (page φ_1 ⊗ ψ_2). p. 179: <λ, x> = λψ(x).
%  - p. 189: e_1 ⊗ e_1 read e_1 ⊗ e_2.
%
% Announced and never established in the folder: the explicit transition
% isomorphisms for A_Q (p. 23, « Je ne vais pas faire l'explicitation ici »);
% an Arf invariant for forms of any rank, even and odd (p. 23, sketched p. 1);
% the canonical isomorphism N_{A/O}(C) ≅ O of p. 32; the « invariants
% d'Hasse-Stiefel-Whitney généralisés » of the cover (p. 42), never reached;
% the modified product ∧^! without « 2 régulier » (p. 55, « exercice
% catégorique universel »); the programme of p. 91 (constraints, twisting,
% étaleness, flatness of X_1 × X_2 -> X_1 * X_2); the quasi-biquadratic
% structure on π^{-1}(ε) without étaleness (pp. 120-121); the fourth
% compatibility of p. 177; the computation of pp. 189-191, which stops.
\documentclass[11pt,a4paper]{article}
\input{../preamble/grothendieck}

\folder{84}
\pages{1}{192}
\dating{[à partir de 1982-vers 1986]}
\watermark{Édition de démonstration\\interprétation personnelle de l'œuvre}
\foldertitle{[Formes quadratiques] : notes manuscrites (s.d.) — lecture modernisée du dossier entier}

\begin{document}

\section*{Résumé}

\begin{resume}
Ces feuillets parlent des formes quadratiques à deux variables,
$ax^{2} + bxy + cy^{2}$, et de leurs compagnes naturelles, les « extensions
quadratiques » : ce qu'on obtient en adjoignant à des nombres une racine
carrée, comme $\mathbf{Z}[\sqrt{-5}]$, ou une racine d'une équation
$x^{2} + bx + c = 0$, comme $\mathbf{Z}[\tfrac{1 + \sqrt{5}}{2}]$. Mais ils en
parlent au-dessus d'une base quelconque --- un anneau comme $\mathbf{Z}$, ou
une famille d'anneaux qui varie d'un point à l'autre d'un espace --- et ils
s'obstinent sur ce qui rend la question difficile : le nombre $2$.

Tant que $2$ est inversible, tout est simple. On « complète le carré » :
$x^{2} + bx + c = (x + \tfrac{b}{2})^{2} - \tfrac{\delta}{4}$, avec
$\delta = b^{2} - 4c$, et une extension quadratique n'est rien d'autre que le
choix d'une racine carrée de $\delta$. Deux extensions se multiplient alors comme
leurs racines : $\sqrt{\delta_1} \cdot \sqrt{\delta_2} = \sqrt{\delta_1\delta_2}$.
Quand $2$ n'est pas inversible, on ne peut plus diviser par $2$, et la racine
carrée ne suffit plus à décrire l'équation. Sur $\mathbf{Z}$, c'est le vieux
fait que le discriminant d'un anneau quadratique est congru à $0$ ou à $1$
modulo $4$, et qu'il faut le savoir pour retrouver l'anneau ; en
caractéristique $2$, c'est pire, car « prendre une racine carrée » n'a plus
rien à voir avec une équation du second degré séparable.

La première partie montre qu'une forme binaire, à un facteur près, est la même
chose qu'une algèbre de matrices $2 \times 2$ commutative --- l'anneau de ses
« homothéties » --- et que donner la forme revient à donner une extension
quadratique et un module de rang un sur elle. C'est, sur une base quelconque,
la correspondance de Gauss et de Dedekind entre formes et idéaux. La deuxième
partie construit le produit de deux extensions quadratiques, sans jamais
diviser par $2$. L'outil est une petite géométrie affine où l'on ne connaît une
origine qu'à un facteur $\chi$ près ; avec $\chi = 2$, elle redonne la
multiplication des racines carrées là où $2$ est inversible, et elle continue de
fonctionner là où il ne l'est pas. La troisième partie reprend cet outil pour
des extensions de modules quelconques.

Le dossier est un atelier. On l'y voit deviner une formule (« Je devine qu'on
a »), la vérifier, s'avouer qu'une hypothèse de départ était fausse (« la
supposition du début est erronée ») et « rectifier le tir », s'arrêter sur un
signe qui revient trois fois sous trois déguisements, et saluer l'objet trouvé
d'un « Plus de doute, c'est lui ! ». La couverture annonce des « invariants
d'Hasse--Stiefel--Whitney généralisés » : le dossier ne les atteint pas.

Les noms sous lesquels chercher la suite sont ceux de la composition de Gauss
sur une base quelconque, de l'algèbre de Clifford paire, de l'invariant d'Arf
et de l'algèbre discriminante, des théories de Kummer et d'Artin--Schreier,
des schémas en groupes d'ordre $2$ de Tate et Oort, et du produit des algèbres
quadratiques.

\keywords{binary quadratic form, quadratic form with values in a line bundle, similarity class, even Clifford algebra, Gauss composition, quadratic algebra, invertible module, norm of an invertible module, Weil restriction, norm-one torus, Hilbert 90, Azumaya algebra, reduced norm, canonical involution, reflection, quadratic cover, discriminant, Stickelberger congruence, Kummer theory, Artin-Schreier theory, Tate-Oort group scheme, contracted product of torsors, extension of modules, product of quadratic algebras, Arf invariant, discriminant algebra, biquadratic algebra, Klein four-group, algebra with standard involution}
\end{resume}

\section*{Le fil du dossier, et les conventions}

\subsection*{Les stations}

\begin{enumerate}
\item[0.] \emph{Feuillets d'ouverture} (pages 1 à 5) : un monoïde gradué par
$\mathbf{Z}/2$ fait de classes de revêtements quadratiques et de couples
$(L, \delta)$ (pages 1 et 2) ; une algèbre triquadratique et ses quatorze
sous-algèbres (page 3) ; les automorphismes d'une algèbre quadratique et ses
éléments de norme $1$ (page 5).
\item[I.] \emph{Formes binaires et revêtements quadratiques} (pages 8 à 41) :
structures de similitude (pages 8 à 11) ; en rang $2$, l'anneau des homothéties
(pages 11 à 18) ; l'algèbre $A_Q$ d'une forme (pages 18 à 23) ; le théorème
récapitulatif (pages 23 à 27) ; la reformulation par une algèbre d'Azumaya
(pages 27 à 34) ; unicité de la conjugaison et symétries (pages 34 à 41).
\item[II.] \emph{$\chi$-scindages et produits contractés} (pages 42 à 121) :
description d'un revêtement quadratique par $(L, P_1, Q)$ et, quand $2$ est
régulier, par $(L, \delta, T_0)$ (pages 43 à 53) ; torseurs sous $\mu_2$ et sous
$\mathbf{Z}/2$ (pages 53 à 55) ; caractéristique $2$ (pages 56 et 57) ; quatre
catégories équivalentes (pages 58 à 61) ; le produit, deviné en coordonnées
(pages 62 à 65), puis construit par la « digression affine » (pages 65 à 79) ;
le cas $\chi = 2$ et le critère des coinvariants (pages 79 à 87) ; la forme
quadratique du produit, $\delta = \delta_1\delta_2$ (pages 88 à 92) ; dualité
entre $A$ et $E$ (pages 93 à 97) ; trace et norme de $A_1 \otimes A_2$ sur
$A_1 * A_2$ (pages 98 à 106) ; algèbres quasi-quadratiques et
quasi-biquadratiques (pages 106 à 121).
\item[III.] \emph{Extensions $\chi$-scindées quelconques} (pages 123 à 191) :
coordonnées et cocycle (pages 123 à 131) ; $\lambda$-rétractions et produit des
morphismes (pages 133 à 139) ; l'opérateur symétrique $F_{11}$ et la condition
$\gamma\chi = 2$ (pages 141 à 157) ; produits poussé et tiré, dualité (pages 159
à 169) ; retour aux algèbres quadratiques : l'accouplement $A \times E \to
\mathcal{O}$ sur le produit (pages 171 à 191).
\end{enumerate}

\subsection*{Conventions, valables pour tout le dossier}

$X$ est un schéma, ou plus généralement un topos localement annelé (c'est le
cadre que pose la page 8), et $\mathcal{O} = \mathcal{O}_X$. Les Modules sont des
$\mathcal{O}$-Modules ; $\Gamma$ désigne les sections, $\check{E} = E^{\vee}$ le dual.
On note
\[
  \underline{\mathrm{Quad}}(E) = \underline{\mathrm{Hom}}(\Gamma^{2}E, \mathcal{O}),
\]
le faisceau des formes quadratiques sur $E$, où $\Gamma^{2}$ est la puissance
divisée ; pour $E$ localement libre de type fini,
$\underline{\mathrm{Quad}}(E) \simeq \underline{\mathrm{Sym}}^{2}(E^{\vee})$.

Une \emph{algèbre quadratique} est une $\mathcal{O}$-algèbre commutative
localement libre de rang $2$ ; on note $\mathrm{Tr}$ et $N$ sa trace et sa norme,
$\sigma(x) = \bar{x} = \mathrm{Tr}(x) - x$ son involution canonique, et
$X' = \operatorname{Spec} A$ le \emph{revêtement quadratique} correspondant. Le
schéma en groupes des unités de $A$ (restriction de Weil de $\mathbb{G}_m$) est
$\mathbb{W}(A)^{*}$, et $\mathcal{U}_A \subset \mathbb{W}(A)^{*}$ le noyau de la
norme\footnote{La page écrit $W(A)^{*}$, $SW(A)$, $\mathcal{U}_A$ ou
$\mathcal{U}(A)$ ; on garde $\mathbb{W}(A)^{*}$ et $\mathcal{U}_A$.}.

À partir de la deuxième partie, $E = \check{A}$, et l'unité de $A$ donne une
extension
\[
  0 \to L \to E \xrightarrow{\ \psi\ } \mathcal{O} \to 0, \qquad \psi(x) = x(1),
\]
où $L \simeq (A/\mathcal{O})^{\vee}$ est inversible ; $P_\lambda = \psi^{-1}(\lambda)$
est un torseur sous $L$. La trace est une section $T$ de $P_2$, la norme une
section $N$ de $\mathrm{Sym}^{2}E$. Le choix d'une section $x_0$ de $P_1$ (un
\emph{scindage}) écrit tout en coordonnées :
\[
  Q(x_0 + u) = u^{2} + bu + c, \qquad b = 2x_0 - T \in \Gamma L, \quad c \in \Gamma L^{\otimes 2}.
\]

\emph{Discriminants.} On note $\delta = b^{2} - 4c = T^{2} - 4N$, section de
$L^{\otimes 2}$ ; c'est le $\delta$ du dossier à partir de la page 43, qu'il
appelle « le discriminant de $N$ changé de signe ». Pour une forme binaire $Q$ de
la première partie, $\mathrm{d}(Q)$ désigne le déterminant de sa forme polaire
(son $\delta(Q)$ de la page 18) ; pour $Q = \alpha\lambda^{2} + \beta\lambda\mu +
\gamma\mu^{2}$, $\mathrm{d}(Q) = 4\alpha\gamma - \beta^{2}$. Ainsi
$\delta = -\mathrm{d}$ partout.

\emph{Produits.} Le produit de deux extensions munies d'une
$\chi$-trivialisation, qu'il note tour à tour $\wedge$ (page 68), $*$, $\circledast$
(page 143) et $\vee$ (troisième partie), est noté $E_1 * E_2$ partout ; le
« produit tiré » de la troisième partie, son $\wedge$ des pages 167 à 169, est
noté $E_1 *^{\varpi} E_2$. Le produit des revêtements quadratiques est
$X_1 * X_2$, celui des algèbres $A_1 * A_2$ (il écrit $X_1 \wedge X_2$,
$X_1 \star X_2$, $A_1 \wedge A_2$).

\emph{Noms réservés.} Dans la première partie, $\Delta$ est le Module inversible
d'une structure de similitude et $\Delta' = \Delta^{\vee}$ ; $\Lambda = \Delta \otimes L$.
L'algèbre d'Azumaya de la section 4 est notée $\mathcal{M}$ (son $M$), comme au
dossier 82, et l'ensemble des $\theta$ tels que $\theta u = \bar{u}\theta$ est noté
$\mathcal{J}$ (son $M$ de la page 39).

\subsection*{Ce que seule une lecture d'ensemble peut dire}

\emph{Les premiers feuillets viennent après.} Les pages 1 et 2 manipulent deux
applications $\alpha : (L, \delta, T_0) \mapsto (L, \delta)$ et
$\beta : (L, \delta) \mapsto (L, 4\delta, 0)$, et constatent que $\beta$ n'est pas
multiplicative. Ce sont exactement les foncteurs $\mathrm{Ass}'$ et
$\mathrm{Oub}'$ des pages 53 et 54, et le constat est celui de la page 54. Les
pages 1 et 2 supposent donc la description des revêtements quadratiques par des
triplets $(L, \delta, T_0)$, qui n'est établie qu'aux pages 47 à 53 ; elles sont
une seconde tentative, à côté de la « catégorie monoïdale mixte » de la page 55,
pour réparer ce défaut.

\emph{Trois régimes, un formalisme.} La deuxième partie traite successivement
trois cas : $2$ inversible, où un revêtement quadratique est une racine carrée
(pages 46 à 48, « la théorie est triviale », page 81) ; $2$ non diviseur de zéro,
où il est un triplet $(L, \delta, T_0)$ (pages 48 à 53) ; caractéristique $2$, où
il est un torseur sous un schéma en groupes d'ordre $2$ (pages 56 et 57). La
« digression affine » (pages 65 à 79) est ce qui les traite tous à la fois, sans
hypothèse sur $2$.

\emph{Un même signe.} « Attention au signe » revient aux pages 17--18, 43, 96 et
133--135 (« signe étrange ! »). Aux pages 17--18 et 43, c'est la convention du
discriminant ($\delta = -\mathrm{d}$, ci-dessus). Aux trois autres endroits, il
nous semble --- ce n'est pas dit sur les pages --- que c'est un seul et même
signe : l'autodualité d'un Module $E$ de rang $2$, $E \simeq
\underline{\mathrm{Hom}}(E, \det E)$, $x \mapsto (y \mapsto x \wedge y)$, induit
$-\mathrm{id}$ sur un sous-module de rang $1$, parce que $x \wedge y = -y \wedge x$. On
le rencontre page 13 dans $\mathcal{O} \to \underline{\mathrm{End}}(E)$, page 96
comme le $-\mathrm{id}_L$ de l'isomorphisme $g$, pages 133 et 135 comme le
$-\mathrm{id}$ de $\Phi$.

\emph{L'ordre des feuillets.} Les pages 35 à 38 ne sont pas dans l'ordre de la
rédaction : la page 38 achève la démonstration de la page 34, la page 37 ouvre
le paragraphe 4.2, que continuent les pages 36 puis 35, et la page 39 suit. On
lit dans cet ordre.

\subsection*{Ce que le dossier annonce sans l'établir}

L'explicitation des isomorphismes de transition de $A_Q$ (page 23, « Je ne vais
pas faire l'explicitation ici ») ; un invariant d'Arf pour les formes de rang
quelconque, pair ou impair (page 23, esquissé page 1) ; l'isomorphisme canonique
$N_{A/\mathcal{O}}(C) \simeq \mathcal{O}$ de la page 32 ; les invariants
d'Hasse--Stiefel--Whitney de la couverture (page 42) ; le produit modifié
$\wedge^{!}$ sans l'hypothèse « $2$ régulier » (page 55) ; le programme de la
page 91 ; la structure biquadratique sur $\pi^{-1}(\varepsilon)$ sans hypothèse
d'étalité (pages 120 et 121) ; la quatrième compatibilité de la page 177 ; le
calcul des pages 189 à 191, qui s'arrête.

\pagerange{1}{5}
\section*{0. Les feuillets d'ouverture (pages 1 à 5)}

\pagerange{1}{2}
\subsection*{Un monoïde gradué par $\mathbf{Z}/2$ (pages 1 et 2)}

Soient $R_0$ l'ensemble des classes de triplets $(L, \delta, T_0)$ décrivant les
revêtements quadratiques (pages 47 à 53 : $L$ inversible, $\delta \in \Gamma
L^{\otimes 2}$, $T_0$ une section de $L$ sur $V(2)$ avec $T_0^{2} \equiv \delta$
modulo $4$), et $R_1$ celui des classes de couples $(L, \delta)$, qui décrivent
les $\mu_2$-torseurs quand $\delta$ est inversible (page 53). Tous deux sont des
monoïdes commutatifs unitaires pour le produit
\[
  (L, \delta, T_0)(L', \delta', T'_0) = (LL', \delta\delta', T_0T'_0), \qquad
  (L, \delta)(L', \delta') = (LL', \delta\delta').
\]
L'oubli $\alpha : R_0 \to R_1$, $(L, \delta, T_0) \mapsto (L, \delta)$, est un
homomorphisme ; l'application $\beta : R_1 \to R_0$, $(L, \delta) \mapsto (L,
4\delta, 0)$, qui fait d'un $\mu_2$-torseur $z^{2} = \delta$ le revêtement
quadratique qu'il est, ne l'est pas\footnote{La page 1 écrit d'abord que
$\alpha$ envoie $(L, \delta, T_0)$ sur « $(L, \Gamma)$ » ; on lit $(L, \delta)$. Elle
qualifie $\beta$ de « multiplicat. », puis ajoute « non compat. avec produits »,
d'une lecture incertaine ; la page 2 tranche : « $R_1 \to R_0$ n'est \emph{pas} un
hom. de monoïde ».}.

Il pose $R_* = R_0 \sqcup R_1$ et définit une loi $\wedge$ compatible à la
graduation par $\mathbf{Z}/2$ :
\[
  x \wedge y = \begin{cases}
    xy & x, y \in R_0, \\
    \beta(xy) & x, y \in R_1, \\
    x\,\alpha(y) & x \in R_1,\ y \in R_0 .
  \end{cases}
\]
L'associativité se ramène, par l'examen des huit cas, à deux identités :
\[
  \text{(A)} \quad \beta(\alpha(x)u) = x\,\beta(u), \qquad
  \text{(B)} \quad \alpha\beta(v) = v\,\eta, \quad \eta = \alpha(\zeta), \ \zeta = \beta(1) = (\mathcal{O}, 4, 0),
\]
pour $x \in R_0$, $u, v \in R_1$. Toutes deux se vérifient sur les formules :
$\beta(\alpha(x)u) = (LL', 4\delta\delta', 0) = x\,\beta(u)$, et
$\alpha\beta(L, \delta) = (L, 4\delta) = (L, \delta)(\mathcal{O}, 4)$. Elles donnent
la mesure exacte du défaut de multiplicativité :
\[
  \beta(x)\,\beta(y) = \zeta\,\beta(xy).
\]
Avec la loi $\wedge$, le produit de $(L, \delta)$ par $(L', \delta', T'_0)$ est
$(LL', \delta\delta')$\footnote{La page 2 écrit, pour ce cas, $(LL', 2\delta\delta')$,
le « $2$ » surchargé et douteux. La définition de la page 1 donne
$\delta\delta'$, et c'est ce qu'exige l'associativité.}.

Dans un encadré, il en tire un homomorphisme. Notons $R_*^{\mathbf{Z}}$ le monoïde
des couples $(n, x)$, $n \in \mathbf{N}$, $x \in R_{n \bmod 2}$, avec
$(n, x)(m, y) = (n + m, x \wedge y)$ ; posons
\[
  w(2\nu, x) = x\,\zeta^{\nu}, \qquad w(2\nu + 1, x) = \beta(x)\,\zeta^{\nu} .
\]
Alors $w : R_*^{\mathbf{Z}} \to R_0$ est un homomorphisme : pour deux degrés
impairs, $w(x \wedge y) = \beta(xy)\zeta^{\nu + \mu + 1} = \beta(x)\beta(y)\zeta^{\nu+\mu}$
par la formule ci-dessus ; pour un degré pair et un impair, c'est (A)\footnote{La
page écrit « $R_* = R \times_{\mathbf{Z}} \mathbf{Z}/2\mathbf{Z}$ » ; on lit le produit
fibré $\mathbf{Z} \times_{\mathbf{Z}/2} R_*$, ce que demandent les formules, qui
dépendent de $n$ et non seulement de sa parité. Le bord de la feuille est déchiré
là où il annonce un second homomorphisme $w' : R_* \to \mathbf{Z} \times R_0$. La
vérification de la multiplicativité est de l'édition.}.

Pour une somme orthogonale $E = \bigoplus_{i \le n}(L_i, \delta_i)$ de $n$ formes de
rang $1$, il compare deux candidats invariants dans $R_0$ :
$w_1(E) = \prod_i \beta(L_i, \delta_i) = (\bigotimes L_i, 4^{n}\prod\delta_i, 0)$, et un
« invariant d'Arf » $(\bigotimes L_i, 2^{\nu}\prod\delta_i, 0)$ avec $\nu$ pair
voisin de $n$, qui en diffère par une puissance de $\zeta$. L'idée lisible est
que, $\zeta$ n'étant pas inversible dans $R_0$ quand $2$ ne l'est pas, l'invariant
d'Arf détermine $w_1$ sans réciproque\footnote{Les exposants ($\nu$, $\nu'$,
$\nu''$) sont en partie incertains, et le $\xi$ de la page désigne tantôt
$(\mathcal{O}_X, 2^{?}, 1)$, tantôt $(\mathcal{O}_X, 4, 0) = \zeta$. La remarque
« $\mathrm{Arf}(E)$ est un invariant plus fin que $c_1(E)$ (si $2$ inversible sur
$X$) » ne peut pas être lue telle quelle : si $2$ est inversible,
$(\mathcal{O}, 4, 0) \simeq (\mathcal{O}, 1, 0)$ par $z \mapsto 2z$, $\zeta$ est la classe
unité, et les deux invariants coïncident. Rien n'établit que l'« invariant d'Arf »
ainsi défini ne dépende pas de la décomposition orthogonale ; c'est une
proposition, non un résultat. La page s'arrête sur « Mais $c_1(E)$ ».}.

\pagerange{3}{3}
\subsection*{Une algèbre triquadratique (page 3)}

Soient $A_i = \mathcal{O}[u_i]/(u_i^{2} + e_i)$, $i = 1, 2, 3$, et
$A = A_1 \otimes A_2 \otimes A_3$, de base les huit monômes en $u_1, u_2, u_3$. Les
involutions canoniques $\sigma_i$ engendrent une action de $\mathbf{F}_2^{3}$ sur
$A$. Quand les $A_i$ sont étales, c'est-à-dire quand les $2e_i$ sont
inversibles\footnote{La page dit « si les $e_i$ inv., i.e.\ les $A_i$ étales ». Le
discriminant de $u^{2} + e$ est $-4e$ : il faut aussi $2$ inversible.}, la
correspondance de Galois pour $\mathbf{F}_2^{3}$ s'écrit en deux tableaux :
\begin{itemize}
\item les sept sous-algèbres quadratiques $A_1, A_2, A_3$, $B_1 = \langle u_2u_3\rangle$,
$B_2 = \langle u_3u_1\rangle$, $B_3 = \langle u_1u_2\rangle$, $B_0 = \langle u_1u_2u_3\rangle$,
fixées chacune par un sous-groupe d'ordre $4$, c'est-à-dire par un hyperplan de
$\mathbf{F}_2^{3}$ ;
\item les sept sous-algèbres de rang $4$, $A'_1 = A_2 \otimes A_3$, \ldots,
$B'_0 = \langle u_2u_3, u_3u_1\rangle$, fixées chacune par une involution, c'est-à-dire
par un élément non nul de $\mathbf{F}_2^{3}$\footnote{La page hésite : le texte
écrit « correspondant aux sept droites vect.\ », la marge corrige en « sept
hyperplans ». Les deux sont justes, l'un pour chaque tableau : les sous-algèbres
quadratiques correspondent aux hyperplans, celles de rang $4$ aux droites, ou aux
éléments non nuls.}.
\end{itemize}
On a $B_1 \simeq A_2 * A_3$, $B_2 \simeq A_3 * A_1$, $B_3 \simeq A_1 * A_2$,
$B_0 \simeq A_1 * A_2 * A_3$ --- le produit des revêtements quadratiques, au sens
de la deuxième partie, qui est ici celui des classes de Kummer :
$(u_2u_3)^{2} = e_2e_3$ --- et $A_i \otimes A'_i \simeq A$, $B_0 \otimes B'_0 \simeq A$. D'où
« $7 + 7 = 14$ situations biquadratiques typiques », chacune une sous-algèbre de
rang $4$ avec les trois sous-algèbres quadratiques qu'elle contient. Ce tableau
est la situation modèle que les pages 110 à 121 axiomatiseront.

\pagerange{5}{5}
\subsection*{Automorphismes et éléments de norme $1$ (page 5)}

Soit $A$ quadratique de base $(1, u)$. Un automorphisme qui induit $-\mathrm{id}$
sur $A/\mathcal{O}$ envoie $u$ sur $\bar{u} + \alpha$ ; conserver trace et norme
impose $2\alpha = 0$ et $\alpha(\mathrm{Tr}\,\bar{u} + \alpha) = 0$, donc $\alpha = 0$ si
$2$ est inversible. La question sans hypothèse sur $2$ est reprise pages 34 et
38.

Les éléments de norme $1$ sont les $u\bar{u}^{-1}$ (théorème 90 de Hilbert pour
l'algèbre $A$), et il le met en diagramme :

\begin{tikzcd}
1 \arrow[r] & \mu_2 \arrow[r] \arrow[d] & \mathbb{G}_m \arrow[r, "(\ )^{2}"] \arrow[d] & \mathbb{G}_m \arrow[r] \arrow[d, no head, "=" description] & 1 \\
1 \arrow[r] & \mathcal{U}_A \arrow[r] & \mathbb{W}(A)^{*} \arrow[r, "N"] & \mathbb{G}_m \arrow[r] & 1
\end{tikzcd}

d'où, par le lemme du serpent, $\mathcal{U}_A/\mu_2 \simeq \mathbb{W}(A)^{*}/\mathbb{G}_m$,
comme faisceaux pour la topologie fppf\footnote{Les deux lignes sont exactes pour
la topologie fppf (l'élévation au carré est fppf-surjective, et la norme aussi).
La page dessine les flèches de bas en haut et biffe en partie la ligne du
quotient ; une flèche porte « cas \ldots », illisible.}. La composée
$\mathcal{U}_A \to \mathcal{U}_A/\mu_2 \simeq \mathbb{W}(A)^{*}/\mathbb{G}_m \xrightarrow{u \mapsto
u\bar{u}^{-1}} \mathcal{U}_A$ est $u \mapsto u^{2}$ ; on la retrouvera page 41.

\pagerange{8}{41}
\section*{I. Formes quadratiques binaires et revêtements quadratiques (pages 8 à 41)}

\pagerange{8}{11}
\subsection*{1. Formes à un scalaire près : structures de similitude (pages 8 à 11)}

Soit $E$ un Module. Le groupe $\mathcal{O}^{*}$ opère sur
$\underline{\mathrm{Quad}}(E)$\footnote{La page écrit
$\underline{\mathrm{Quad}}(E) \simeq \underline{\mathrm{Hom}}(\Gamma^{2}(E), E)$ ; le but est
$\mathcal{O}$. L'homomorphisme canonique $\underline{\mathrm{Sym}}^{2}(E^{\vee}) \to
\underline{\mathrm{Quad}}(E)$, $fg \mapsto (x \mapsto f(x)g(x))$, est un isomorphisme
pour $E$ localement libre de type fini, comme elle le dit.}, et une section du
faisceau quotient
\[
  \underline{\mathrm{Sim}}(E) = \underline{\mathrm{Quad}}(E)/\mathcal{O}^{*}
\]
est une \emph{structure de similitude} sur $E$ : localement une forme $Q$, deux
formes $Q$ et $Q' = aQ$, $a$ inversible, définissant la même. Elle est
\emph{fidèle} si $a \mapsto aQ$, $\mathcal{O} \to \underline{\mathrm{Quad}}(E)$, est
injective --- le scalaire $a$ est alors unique, et la condition ne dépend que
de la structure.

Une structure de similitude $\Sigma$ est déterminée par le sous-Module $\Delta_\Sigma
\subset \underline{\mathrm{Quad}}(E)$ qu'elle engendre, et $\Sigma \mapsto \Delta_\Sigma$
est une bijection sur les sous-Modules localement monogènes\footnote{C'est ici
que sert l'hypothèse d'un topos \emph{localement} annelé : deux générateurs d'un
module monogène $\mathcal{O}_x Q$ sur un anneau local diffèrent par une unité. La
page ne le dit pas.}. Les structures fidèles correspondent aux sous-Modules
inversibles ; les structures \emph{strictement} (ou universellement) fidèles,
celles pour lesquelles $\Delta \to \underline{\mathrm{Quad}}(E)$ reste injectif sur
chaque fibre résiduelle, aux sous-fibrés en droites de
$\underline{\mathrm{Quad}}(E)$, c'est-à-dire aux sections du fibré projectif des
droites de $\underline{\mathrm{Quad}}(E)$\footnote{Pages 9 et 10, on suppose $E$
localement libre de type fini. La page 9 écrit
$\underline{\mathrm{Quad}}(E) \simeq \underline{\mathrm{Sym}}^{2}(E)$, la page 8
$\underline{\mathrm{Sym}}^{2}(E^{\vee})$ ; c'est la seconde. Le paragraphe 1.3 est
très surchargé, et un « NB » interlinéaire y est en partie illisible.}.

\emph{La forme canonique} (1.4). Soient $\Delta$ inversible et $i : \Delta \to
\underline{\mathrm{Quad}}(E)$ un homomorphisme quelconque. La forme
$Q_E : E \to \underline{\mathrm{Quad}}(E)^{\vee}$, $\langle f, Q_E(x)\rangle = f(x)$, est
quadratique, et sa composée avec ${}^{t}i$ est une forme
\[
  Q_i : E \longrightarrow \Delta^{\vee} = \Delta'
\]
\emph{à valeurs dans un Module inversible}. Pour une structure fidèle, $(\Delta, i)$
est déterminé par $\Sigma$ à isomorphisme unique près, d'où la notation
$Q_\Sigma$. Si l'on choisit une base de $\Delta$, $Q_\Sigma$ redevient une forme
ordinaire. D'où (1.5, en marge) : une structure de similitude strictement fidèle
équivaut à la donnée d'un Module inversible $\Delta'$ et d'une forme quadratique
$Q : E \to \Delta'$ qui n'est nulle sur aucune fibre résiduelle\footnote{Une forme
quadratique non nulle sur un espace vectoriel prend une valeur non nulle, sur
tout corps, $\mathbf{F}_2$ compris ; « $Q(E)$ engendre $\Delta'$ » et « $Q$ non nulle
sur toute fibre » sont donc équivalents. La marge écrit « fidèle », surchargé de
« strictement ».}. C'est la notion de forme quadratique à valeurs dans un fibré
en droites, sous laquelle on la trouve aujourd'hui.

\pagerange{11}{16}
\subsection*{2. En rang $2$ : l'anneau des homothéties (pages 11 à 16)}

Désormais $E$ est localement libre de rang $2$, et $L = \det E = \bigwedge^{2}E$.
L'accouplement $(x, y) \mapsto x \wedge y$ donne $E \simeq \check{E} \otimes L$, d'où
\[
  \underline{\mathrm{Quad}}(E) \simeq \underline{\mathrm{Quad}}(\check{E}) \otimes L^{\otimes -2}
\]
et un isomorphisme canonique $\underline{\mathrm{Sim}}(E) \simeq
\underline{\mathrm{Sim}}(\check{E})$, involutif, qui respecte fidélité, stricte
fidélité et non-dégénérescence (2.1).

\emph{La clef} (2.2). La suite exacte $0 \to \bigwedge^{2}E \to E \otimes E \to
\mathrm{Sym}^{2}E \to 0$, $x \wedge y \mapsto x \otimes y - y \otimes x$, tensorisée par
$L^{\vee}$, s'écrit, puisque $E \otimes L^{\vee} \simeq \check{E}$ et $E \otimes
\check{E} \simeq \underline{\mathrm{End}}(E)$ :

\begin{tikzcd}[column sep=small]
0 \arrow[r] & L \otimes L^{\vee} \arrow[r] \arrow[d, "\wr"] & E \otimes E \otimes L^{\vee} \arrow[r] \arrow[d, "\wr"] & \mathrm{Sym}^{2}E \otimes L^{\vee} \arrow[r] \arrow[d, "\wr"] & 0 \\
0 \arrow[r] & \mathcal{O} \arrow[r] & \underline{\mathrm{End}}(E) \arrow[r] & \underline{\mathrm{End}}(E)/\mathcal{O} \arrow[r] & 0
\end{tikzcd}

Le sous-module $L \otimes L^{\vee}$ va sur les homothéties $\mathcal{O}\cdot 1_E$\footnote{La
page écrit la flèche de gauche $\mathrm{id}_{\mathcal{O}}$. Avec les conventions
écrites ($x \mapsto (y \mapsto x \wedge y)$, $x \otimes f \mapsto f(\cdot)\,x$), un
calcul en base donne $1 \mapsto -1_E$. Comme sous-modules, rien ne change ; c'est
le signe du « Ce que seule une lecture d'ensemble peut dire ».}, et l'on obtient
\[
  \underline{\mathrm{End}}(E)/\mathcal{O}\cdot 1_E \simeq \mathrm{Sym}^{2}E \otimes L^{\vee}
  \simeq \underline{\mathrm{Quad}}(E) \otimes L .
\]
\emph{Une structure de similitude sur $E$ est donc la même chose qu'un sous-Module
localement monogène de $\underline{\mathrm{End}}(E)/\mathcal{O}$}, c'est-à-dire qu'un
sous-Module $A \subset \underline{\mathrm{End}}(E)$ contenant $\mathcal{O}\cdot 1_E$ avec
$A/\mathcal{O}$ localement monogène ; les structures strictement fidèles
correspondent aux sous-fibrés de rang $2$ contenant $\mathcal{O}\cdot 1_E$.

\textbf{Lemme 1.} Un tel $A$ est localement de la forme $\mathcal{O} + \mathcal{O}u$ ;
c'est une sous-algèbre commutative de $\underline{\mathrm{End}}(E)$, finie sur
$\mathcal{O}$, puisque $u^{2} - \mathrm{Tr}(u)u + \det(u) = 0$ (Cayley--Hamilton).
On l'appelle \emph{l'anneau des homothéties} (directes) de $\Sigma$, et on le note
$A_\Sigma$\footnote{La condition a') du lemme est illisible. La marge renvoie la
finitude à l'équation $u^{2} - \mathrm{Tr}(u)u + N(u) = 0$.}.

\textbf{Lemme 2.} Pour $\Sigma$ de la forme $A = \mathcal{O} + \mathcal{O}u$, sont
équivalentes : $\Sigma$ strictement fidèle ; $A$ localement libre de rang $2$ et
$A(s) \to \underline{\mathrm{End}}(E(s))$ injectif en tout point ; $u(s)$ n'est
scalaire en aucun point ; $E$ est un $A$-Module inversible ; $E$ est un $A$-Module
localement monogène ; $\underline{\mathrm{End}}(E)$ est un $A$-Module localement
libre de rang $2$\footnote{Le passage de « $u(s)$ non scalaire » à « $E$ localement
monogène sur $A$ » est le fait qu'une matrice $2 \times 2$ non scalaire admet un
vecteur cyclique, suivi de Nakayama ; la page ne détaille pas. Le dossier 82,
pages 36 à 40, étudie le même commutant $\mathcal{O} + \mathcal{O}u$ d'un
endomorphisme nulle part scalaire, côté $\mathrm{PGL}_2$.}.

\textbf{Corollaire 1.} Alors $\Delta \simeq (A/\mathcal{O}) \otimes L^{\vee}$, et la forme
de définition est
\[
  Q_\Sigma : E \longrightarrow \Delta' \simeq (A/\mathcal{O})^{\vee} \otimes L.
\]\footnote{La
  page écrit $\Delta' \simeq A/\mathcal{O} \otimes L^{\vee}$, ce qui est $\Delta$ et non
  son dual.}
Pour une forme $Q : E \to \Delta'$ et une section $x$ de $E$, sont équivalentes :
$a \mapsto ax$, $A \to E$, est surjectif ; c'est un isomorphisme ; $Q(x)$ est une
base de $\Delta'$ ; en tout point, $x(s)$ n'est pas vecteur propre de $u(s)$.

\pagerange{17}{18}
\subsection*{La forme $Q_u$ et le discriminant (pages 17 et 18)}

L'image d'un endomorphisme $u$ dans $\underline{\mathrm{End}}(E)/\mathcal{O} \simeq
\underline{\mathrm{Quad}}(E, L)$ est la forme
\[
  Q_u(x) = x \wedge u(x),
\]
nulle si et seulement si $u$ est une homothétie. En coordonnées, pour
$u = \begin{pmatrix} a & b \\ c & d \end{pmatrix}$,
$Q_u(\lambda, \mu) = c\lambda^{2} + (d - a)\lambda\mu - b\mu^{2}$, et les $u$ tels que
$Q_u = \alpha\lambda^{2} + \beta\lambda\mu + \gamma\mu^{2}$ sont les
\[
  u_a = \begin{pmatrix} a & -\gamma \\ \alpha & a + \beta \end{pmatrix}, \qquad a \text{ quelconque}
\]
--- parmi lesquels $u_0$, de trace $\beta$ et de déterminant $\alpha\gamma$\footnote{La
formule encadrée de la page écrit $Q_u(x) = u(x) \wedge x$, « il faut vérifier le
signe ! » ; ses coordonnées, et les matrices $u_a$, suivent $x \wedge u(x)$, qu'on
retient. Le signe devant $b\mu^{2}$ est couvert d'une tache ; le calcul donne $-$.}.

Pour une forme $Q : E \to \Delta'$, la forme polaire $\varphi_Q$ donne
$E \to \check{E} \otimes \Delta'$, dont le déterminant est une section
\[
  \mathrm{d}(Q) \in \Gamma\bigl((\Delta' \otimes L^{\vee})^{\otimes 2}\bigr),
\]
section de $\mathcal{O}$ quand $\Delta' = L$\footnote{Une ligne antérieure écrit
$\Gamma(\Delta'^{\otimes 2} \otimes L^{\otimes 2})$ ; le calcul qui suit donne
$L^{\vee \otimes 2}$.}. Et l'on a
\[
  \boxed{\ (\mathrm{Tr}\,u)^{2} - 4\det u = -\,\mathrm{d}(Q_u)\ }
\]
--- « attention au signe » : avec $\mathrm{d}(Q) = 4\alpha\gamma - \beta^{2}$, les deux
membres valent $(a - d)^{2} + 4bc$. Dans les conventions de ce texte, le
discriminant $\delta$ de l'algèbre engendrée par $u$ est celui, $\beta^{2} -
4\alpha\gamma$, de la forme.

\pagerange{18}{23}
\subsection*{L'algèbre $A_Q$ d'une forme (pages 18 à 23)}

Soit $Q : E \to \Delta'$ quelconque, nulle ou non. Elle équivaut à
$u_Q : \Lambda = \Delta \otimes L \to \underline{\mathrm{End}}(E)/\mathcal{O}$, et l'on
définit $A_Q$ comme l'image réciproque de l'extension des endomorphismes :

\begin{tikzcd}
0 \arrow[r] & \mathcal{O} \arrow[r] \arrow[d, no head, "=" description] & A_Q \arrow[r] \arrow[d, "u^{!}_Q"] & \Lambda \arrow[r] \arrow[d, "u_Q"] & 0 \\
0 \arrow[r] & \mathcal{O} \arrow[r] & \underline{\mathrm{End}}(E) \arrow[r] & \underline{\mathrm{End}}(E)/\mathcal{O} \arrow[r] & 0
\end{tikzcd}

$A_Q$ est localement libre de rang $2$, muni de la section $1_Q$ et de la forme
$N_Q = \det \circ\, u^{!}_Q$, qui vaut $1$ en $1_Q$. Or un Module localement libre
de rang $2$ muni d'une section $e$ et d'une forme quadratique $N$ avec $N(e) = 1$
porte une unique structure d'algèbre commutative d'unité $e$ et de norme $N$ :
$x^{2} = \varphi_N(x, e)\,x - N(x)\,e$ (c'est l'équivalence (1') $\Leftrightarrow$ (2) de
la page 58). $A_Q$ est donc une algèbre quadratique, et $u^{!}_Q$ un homomorphisme
d'algèbres qui respecte trace, norme et involution canonique
$v \mapsto \bar{v} = \mathrm{Tr}(v) - v$ de $\underline{\mathrm{End}}(E)$.

La construction est fonctorielle pour les isomorphismes et commute aux
changements de base ; $A_Q \simeq A_{\lambda Q}$ pour $\lambda$ inversible, non pour
$\lambda = 0$. Sont équivalents : $u^{!}_Q$ universellement injectif ; $Q$ non nulle
sur toute fibre résiduelle ; $\Delta \to \underline{\mathrm{Quad}}(E)$ universellement
injectif ; $\Sigma_Q$ strictement fidèle --- et alors $A_Q = A_\Sigma$\footnote{La
page écrit « structure fidèle » au point d) ; c'est la stricte fidélité que
demande l'équivalence, et que donne le lemme 2 qu'elle invoque.}. Pour $Q = 0$,
$A_Q = \mathcal{O} \oplus \Lambda$ avec $\Lambda^{2} = 0$, l'algèbre des nombres duaux,
et $u^{!}_Q$ tue $\Lambda$\footnote{Le nom de l'algèbre n'est pas lu à la page 22 ;
la phrase qui précède le donne.}.

\emph{En coordonnées} (2.6). Pour $Q = \alpha\lambda^{2} + \beta\lambda\mu + \gamma\mu^{2}$,
\[
  A_Q = \mathcal{O}[u]/(u^{2} - \beta u + \alpha\gamma), \qquad
  u \longmapsto u_0 = \begin{pmatrix} 0 & -\gamma \\ \alpha & \beta \end{pmatrix}.
\]
C'est l'\emph{algèbre de Clifford paire} de $Q$ : dans $C(Q)$, $w = e_1e_2$ vérifie
$w^{2} = e_1(\beta - e_1e_2)e_2 = \beta w - \alpha\gamma$\footnote{L'identification
est de l'édition ; le mot n'est pas sur les pages. Pour les formes à valeurs dans
un fibré en droites, l'algèbre de Clifford paire a été définie par Bichsel et
Knus (1994) ; la correspondance entre formes binaires et couples (algèbre
quadratique, module) sur une base quelconque est chez Kneser (\emph{J. Number
Theory}, 1982) et M.~M. Wood (\emph{Adv. Math.}, 2011). Références de l'édition,
citées de mémoire ; rien n'indique que l'auteur connût ces travaux, dont le
premier lui est contemporain.}. Les isomorphismes de transition quand on change
de bases de $E$ et de $\Delta$ ne sont pas explicités : « Je ne vais pas faire
l'explicitation ici ».

Le revêtement quadratique associé à $A_Q$ « mérite de s'appeler \emph{l'invariant
d'Arf} de la forme binaire $Q$ ». Il annonce un invariant d'Arf pour les formes
scalaires de rang quelconque, en distinguant rang pair et rang impair ; le dossier
ne le construit pas\footnote{Deux paragraphes portent le numéro 2.6 (pages 22 et
23). Le nom est conforme à l'usage actuel : pour une forme régulière de rang pair,
le centre de l'algèbre de Clifford paire est l'algèbre discriminante, dont la
classe est l'invariant d'Arf, et en rang $2$ l'algèbre de Clifford paire est
elle-même commutative.}.

\pagerange{23}{27}
\subsection*{3. Le théorème récapitulatif (pages 23 à 27)}

\textbf{Théorème.} Les catégories fibrées en groupoïdes suivantes, sur la catégorie
des schémas, sont équivalentes :
\begin{enumerate}
\item[(1)] les couples $(E, \Delta)$, $E$ localement libre de rang $2$, $\Delta \subset
\underline{\mathrm{Quad}}(E)$ un sous-fibré en droites (structures de similitude
strictement fidèles) ;
\item[(1')] les triplets $(E, \Delta', Q)$, $Q : E \to \Delta'$ nulle sur aucune fibre
résiduelle ;
\item[(2)] les couples $(A, E)$, $A$ algèbre quadratique, $E$ un $A$-Module inversible ;
\item[(2')] les couples $(X', \Lambda')$, $X'$ revêtement quadratique de $X$, $\Lambda'$
inversible sur $X'$.
\end{enumerate}

(1) $\Leftrightarrow$ (1') est la section 1 ; (2) $\Leftrightarrow$ (2') est le dictionnaire des
spectres ; (1) $\Rightarrow$ (2) est la section 2. Dans l'autre sens, pour $p : X' \to X$,
\[
  E = p_{*}\Lambda', \qquad \Delta' = N_{X'/X}(\Lambda'), \qquad Q : p_{*}\Lambda' \to N_{X'/X}(\Lambda') \ \text{l'application norme},
\]
la norme d'un Module inversible le long d'un morphisme fini localement
libre\footnote{Sur la norme le long d'un tel morphisme, voir le dossier 98, pages
2 à 8 ; aujourd'hui EGA~IV, 21.5, ou SGA~4, XVII, 6.3. Le théorème est, sur une
base quelconque, la correspondance de Gauss et Dedekind entre formes binaires
primitives et modules inversibles sur les ordres quadratiques ; voir la note sur
Kneser et Wood ci-dessus.}.

\textbf{Corollaire 1.} Ces catégories équivalent aussi à celle des couples $(A, P)$,
$P$ un torseur sous $\mathbb{W}(A)^{*}$, schéma en groupes lisse de dimension
relative $2$. En termes de $(E, Q)$, $P$ est le schéma des sections locales $x$ de
$E$ telles que $Q(x)$ soit une base de $\Delta'$ ; inversement $\Delta'$ est le fibré
en droites déduit de $P$ par la norme $\mathbb{W}(A)^{*} \to \mathbb{G}_m$. Il revient au
même de prendre des torseurs pour la topologie de Zariski ou fpqc, un module
inversible sur une algèbre finie sur un anneau local étant libre.

\textbf{Corollaire 2.} Les couples $(E, Q)$, $Q : E \to \mathcal{O}$ \emph{à valeurs
scalaires} et nulle sur aucune fibre, équivalent aux couples $(A, \Pi)$, $\Pi$ un
torseur sous $\mathcal{U}_A$ ; $\Pi$ est le schéma des $x$ tels que $Q(x) = 1$. Ici
la topologie de Zariski ne suffit plus\footnote{La marge dit que fpqc et étale
« reviennent au même ». C'est vrai quand $\mathcal{U}_A$ est lisse, par exemple si $A$
est étale ou si $2$ est inversible ; pour $A = \mathcal{O}[\varepsilon]$, $\varepsilon^{2} =
0$, en caractéristique $2$, $\mathcal{U}_A$ n'est pas lisse et il faut la topologie
fppf.}.

\emph{Les racines de la forme.} Le revêtement quadratique de $Q$ est le
sous-schéma $Q = 0$ du fibré projectif des droites de $E$ : à une section de
$\operatorname{Spec} A$, c'est-à-dire à un homomorphisme $A \to \mathcal{O}$ de noyau $J$,
répond la droite $D = J\cdot E$, sur laquelle $Q$ est nulle, et toute droite
isotrope est de cette forme. (On peut supposer $E = A$ et $Q$ la norme, et
c'est immédiat.)

\pagerange{27}{34}
\subsection*{4. Le point de vue de l'algèbre d'Azumaya (pages 27 à 34)}

La notion de structure de similitude, et ses variantes, « ne dépend que de $E$ à
tensorisation près par un Module inversible », donc seulement de l'algèbre
d'Azumaya $\mathcal{M} = \underline{\mathrm{End}}(E)$, de rang $4$\footnote{La page écrit
« de rang $4$ sur $E$ » ; c'est sur $\mathcal{O}$.}. Pour une algèbre d'Azumaya
$\mathcal{M}$ de rang $4$ quelconque (une algèbre de quaternions), une structure de
similitude est un sous-Module localement monogène de
$\mathcal{M}/\mathcal{O}\cdot 1$, qui « joue le rôle du Module des formes
quadratiques » ; dans le cas strictement fidèle, le revêtement quadratique associé
est l'intersection de $\mathbb{P}(A) \subset \mathbb{P}(\mathcal{M})$ avec la quadrique
de la norme réduite, $\mathrm{Nrd}_{\mathcal{M}}(x) = 0$ --- tautologiquement, puisque
la norme réduite induit sur $A$ sa norme\footnote{Le dossier 82, pages 46 à 55,
développe la même géométrie : algèbre de quaternions, norme réduite, involution
canonique, et la conique des éléments de trace et de norme nulles.}.

\emph{Que devient le $A$-Module $E$ ?} Il n'est plus défini qu'à un facteur près
venu de la base, $\Lambda \mapsto \Lambda \otimes p^{*}\Theta$ sur $X'$. Mais
$\Lambda \otimes \sigma^{*}\Lambda^{-1}$ ne change pas, et sa norme est triviale. On
s'attend donc à trouver un $A$-Module canonique, ou, ce qui revient au même, une
réduction à $\mathcal{U}_A$ du groupe structural $\mathbb{W}(A)^{*}$.

\emph{Une hypothèse fausse} (pages 29 à 31). Le candidat évident est
$\mathcal{M}/A$, inversible sur $A$ dans le cas strictement fidèle. Pour
$\mathcal{M} = \underline{\mathrm{End}}(E)$, il cherche un isomorphisme canonique,
à scalaire près, $E \simeq \underline{\mathrm{End}}(E)/A$. Le calcul des déterminants
sur $A$ des deux structures, gauche et droite, de $\underline{\mathrm{End}}(E) \simeq E
\otimes \check{E}$ donne
\[
  \underline{\mathrm{End}}(E)/A \simeq \det\nolimits_A \underline{\mathrm{End}}(E) \simeq
  E^{\otimes_A 2} \otimes \check{L} \simeq \check{E}^{\otimes_A 2} \otimes L .
\]
Un isomorphisme $E \simeq \underline{\mathrm{End}}(E)/A$ reviendrait à une base de
$E \otimes \check{L}$ sur $A$, c'est-à-dire à une droite canonique dans $E$ ; il n'y
en a pas. « Donc la supposition du début est erronée. »\footnote{La page 29 se
termine sur un passage barré (une identification de $\check{\mathbb{P}}(E)$ avec les
droites de norme nulle des endomorphismes de trace nulle) ; la page 30 note que
le membre de droite cherché est une conique $C$ dans $\mathbb{P}(\mathcal{M}^{+})$,
qui n'est pas en général le projectivisé d'un fibré de rang $2$. La première
formule de la page 31 porte une surcharge ; on suit la page 30.}

\emph{En rectifiant le tir} (pages 32 et 33). Il pose, localement sur une base
$(1, u)$ de $A$,
\[
  \varphi_u : \mathcal{M} \to \mathcal{M}, \qquad \theta \longmapsto u\theta - \theta\bar{u},
\]
$A$-linéaire pour la structure gauche, qui ne dépend que de l'image de $u$ dans
$A/\mathcal{O}$ ; d'où une suite exacte de $A$-Modules à gauche
$\mathcal{M} \otimes \check{\Lambda} \to \mathcal{M} \to C \to 0$, $C = \operatorname{Coker}
\varphi_u$, et le Module cherché $E \otimes_A \check{E}^{\otimes_A -1} \simeq C$\footnote{La
page écrit $\mathcal{M}/\mathcal{M}^{\natural}$, avec un signe lu comme un bécarre, et
note le conoyau $\Delta$ ; on l'appelle $C$ pour ne pas le confondre avec le
$\Delta$ de la section 1.}. Il faudrait maintenant un isomorphisme canonique
$N_{A/\mathcal{O}}(C) \simeq \mathcal{O}$ ; il n'est pas construit. L'autodualité
$\langle\theta, \theta'\rangle = \mathrm{Trd}(\theta\bar{\theta}')$ de $\mathcal{M}$
donne ${}^{t}\varphi_u = \varphi_{\bar{u}} = -\varphi_u$, donc le dual de
$\operatorname{Coker}\varphi_u$ est $\operatorname{Ker}\varphi_u$, et à un Module inversible de
la base près, $C$ est le dual sur $A$ de $\operatorname{Ker}\varphi_u$.

Or les sections inversibles de ce noyau ont un sens clair :
\[
  \{\theta \in \mathcal{M}^{*} \mid \theta u = \bar{u}\theta\}
  = \{\theta \text{ qui normalise } A \text{ et y induit } \sigma\},
\]
pseudo-torseur sous le centralisateur $A^{*}$ de $A$, et torseur puisque, dans
$\mathbb{M}_2(\mathcal{O})$, $u$ et $\bar{u}$, de même trace et même norme et nulle
part scalaires, sont localement conjugués.

\pagerange{34}{38}
\subsection*{L'unicité de la conjugaison (pages 34 et 38)}

\textbf{Proposition} (le corollaire de la page 35, précisé). Soit $A$ une algèbre
quadratique, étale en tout point de
$X_0 = V(2\cdot 1_X)$. Tout automorphisme d'algèbre $\theta$ de $A$ induisant
$-\mathrm{id}$ sur $A/\mathcal{O}$ coïncide, au voisinage de chaque point, avec
l'involution canonique $\sigma$ ou avec l'identité, et le second cas ne se produit
qu'au voisinage de points où $2 = 0$. En particulier $\sigma$ est le seul tel
automorphisme qui ne soit l'identité au voisinage d'aucun point.

\emph{Démonstration} (pages 34 et 38). Localement $A = \mathcal{O}[u]/(u^{2} + bu + c)$
et $\theta(u) = \bar{u} + \alpha$ ; trace et norme conservées donnent $2\alpha = 0$ et
$\alpha(\alpha - b) = 0$. Hors de $X_0$, $\alpha = 0$. En $s \in X_0$, $b_s$ est
inversible (étalité). Si $\alpha_s$ est inversible, $\alpha = b$ près de $s$, donc
$2b = 0$, donc $2 = 0$ près de $s$, et $\theta(u) = -b - u + b = u$ : $\theta$ est
l'identité près de $s$. Sinon $\alpha_s \in \mathfrak{m}_s$, $\alpha_s - b_s$ est
inversible, et $\alpha = 0$ près de $s$\footnote{La page 34 affirme d'abord, dans un
passage barré, l'alternative $\alpha(s) = 0$ ou $\alpha(s) = b(s)$ ; la suite,
reliée par un trait, la remplace, et la page 38 l'achève (« cqfd »). Le corollaire
de la page 35 énonce l'unicité sans la restriction ; or en caractéristique $2$
l'identité induit aussi $-\mathrm{id} = \mathrm{id}$ sur $A/\mathcal{O}$, et la
proposition de la même page ajoute bien une condition, « l'identité au voisinage
\ldots », que sa marge dit nécessaire seulement en certains points. On énonce ce
que la démonstration établit. Le même phénomène, $gug^{-1} = u + \mathrm{Tr}\,u$ en
caractéristique $2$, est au dossier 82, page 40.}.

Sans étalité, l'unicité tombe : pour $A = \mathcal{O}[u]/(u^{2})$, les
automorphismes induisant $-\mathrm{id}$ sur $A/\mathcal{O}$ sont les $u \mapsto -u +
\alpha$ avec $2\alpha = 0$, $\alpha^{2} = 0$, et en caractéristique $2$ les
$u \mapsto u + \alpha$, $\alpha^{2} = 0$ --- nombreux sur une base non réduite,
réduits à l'identité sur une base réduite.

\pagerange{35}{41}
\subsection*{4.2. Symétries (pages 37, 36, 35, 39 à 41)}

Soient $(E, \Delta', Q)$ strictement fidèle et $D \subset E$ un sous-fibré en droites
\emph{totalement non isotrope} : $Q|D$ non nulle sur chaque fibre. Quitte à
tensoriser par $D^{-1}$, on suppose $D = \mathcal{O}e$, $Q$ scalaire et $Q(e) = 1$.
Alors $f : A \to E$, $a \mapsto ae$, est un isomorphisme de modules quadratiques,
puisque $Q(ax) = N(a)\,Q(x)$ ; il transporte la structure d'algèbre de $A$ en
l'unique structure d'unité $e$ et de norme $Q$. On \emph{définit} la symétrie
$\sigma_D$ en transportant l'involution canonique :
\[
  \sigma_D(ae) = \bar{a}\,e, \qquad \text{c'est-à-dire} \qquad \sigma_D(x) = \varphi_Q(e, x)\,e - x .
\]
Elle induit $\mathrm{id}$ sur $D$, $-\mathrm{id}$ sur $E/D$, et $\det\sigma_D = -1$. Si
$2$ est inversible, $\sigma_D(x) = 2(e\cdot x)\,e - x$ : c'est la symétrie
orthogonale d'axe $D$ du plan, l'opposée de la réflexion d'hyperplan
$e^{\perp}$\footnote{En général $D^{\perp}$ n'est pas un supplémentaire de $D$, d'où
la définition par $E/D$. La page écrit « symétrie p.r.\ à la \emph{droite} $D$ ».}.

\textbf{Proposition} (page 35). Si $A$ est étale en tout point de $X_0$, $\sigma_D$ est
l'unique automorphisme de la structure de similitude qui induise $\mathrm{id}$ sur
$D$, $-\mathrm{id}$ sur $E/D$, et ne soit l'identité au voisinage d'aucun point.
C'est la proposition précédente, transportée par $f$.

\textbf{Lemme} (page 39). $\sigma_D$ normalise $A \subset \underline{\mathrm{End}}(E)$ et y
induit l'involution canonique : $\sigma_D\,u\,\sigma_D^{-1} = \bar{u}$. Pour $E = A$,
$D = \mathcal{O}\cdot 1$, c'est $\overline{u\bar{x}} = \bar{u}x$.

\textbf{Proposition} (pages 39 et 40). Soient
\[
  \mathcal{J} = \{\theta \in \underline{\mathrm{End}}(E) \mid \theta u = \bar{u}\theta \ \text{pour } u \in A\}, \qquad
  \Pi = \{\theta \in \mathcal{J} \mid \det\theta = -1\}.
\]
\begin{enumerate}
\item[a)] $\mathcal{J}$ est un sous-fibré de rang $2$ de $\underline{\mathrm{End}}(E)$,
sous-$A$-Module à gauche et à droite, les deux structures liées par $u\theta =
\theta\bar{u}$ ; c'est un $A$-Module inversible.
\item[b)] $\Pi$ est un torseur sous $\mathcal{U}_A$ pour la multiplication, localement
trivial pour Zariski, puisqu'il contient les $\sigma_D$\footnote{La page 37 définit
le même torseur par la condition « Norme $1$ » ; les pages 39 et 40 prennent
$\det\theta = -1$, qui est la bonne normalisation : elle contient les symétries.
Un essai de démonstration est barré page 40.}.
\item[c)] Soit $P^{*} = \mathbb{P}(E) - X'$ le schéma des droites totalement non
isotropes de $E$\footnote{La page dit « de $A$ ».}. $\mathbb{W}(A)^{*}$ y opère, à
travers $\mathbb{W}(A)^{*}/\mathbb{G}_m$, et $D \mapsto \sigma_D$ est une application
$f : P^{*} \to \Pi$ telle que
\[
  f(u\cdot D) = u\bar{u}^{-1} f(D) \qquad (u \in \mathbb{W}(A)^{*}),
\]
compatible avec l'homomorphisme $u \mapsto u\bar{u}^{-1}$ de la page 5 ; sur
$\mathcal{U}_A$, celui-ci est $u \mapsto u^{2}$ (page 41). En effet $\sigma_{uD} =
u\sigma_Du^{-1} = u\bar{u}^{-1}\sigma_D$ par le lemme.
\end{enumerate}
Le module $\mathcal{J}$, inversible sur $A$, et son torseur $\Pi$ sous $\mathcal{U}_A$,
sont la réponse à la question de la page 28 que les pages 32 et 33 avaient
approchée : $\mathcal{J}$ est, au signe près de l'involution, le noyau de
$\varphi_u$\footnote{Ce rapprochement est de l'édition ; les pages ne le font
pas.}. La page 41 s'arrête sur l'équivariance.

\pagerange{42}{121}
\section*{II. $\chi$-scindages, produits contractés, revêtements quadratiques (pages 42 à 121)}

La feuille de couverture, sur papier à musique, numérotée « (79) », porte le
titre : « $\chi$-scindages, produits contractés d'extensions $\chi$-scindées,
produits contractés de revêtements quadratiques, invariants
d'Hasse--Stiefel--Whitney généralisés, etc. »\footnote{La transcription lit
« $X$-scindages » ; c'est la lettre $\chi$ des pages 65 et suivantes. Les invariants
d'Hasse--Stiefel--Whitney n'apparaissent nulle part dans le dossier.}

\pagerange{43}{47}
\subsection*{Un revêtement quadratique comme fonction polynôme (pages 43 à 47)}

Avec les conventions posées plus haut, $N(1) = 1$ et $T(1) = 2$ ; $\mathrm{Sym}^{2}E$
a une suite de composition de quotients $\mathcal{O}$, $L$, $L^{\otimes 2}$, et
\[
  \boxed{\ T^{2} - 4N = \delta\ }, \qquad \delta \in \Gamma L^{\otimes 2}.
\]
Dans le cas décomposé $A = \mathcal{O} \times \mathcal{O}$, les deux projections
$x', x''$ donnent $T = x' + x''$, $N = x'x''$, $\delta = (x' - x'')^{2}$\footnote{Le
signe de la seconde formule encadrée est couvert d'une tache ; on le donne tel que
le calcul l'exige. La marge, « attention au signe », rappelle que $\delta$ est le
discriminant de $N$ changé de signe. Elle note aussi $\varphi_N(x, y) =
\mathrm{Tr}(x)\mathrm{Tr}(y) - \mathrm{Tr}(xy)$.}.

Ce cas suggère la fonction $x \mapsto (x - x')(x - x'')$ sur $P_1$, c'est-à-dire
$x \mapsto x^{2} - Tx + N$, et la forme quadratique sur $E$
\[
  Q(x) = x^{2} - \psi(x)\,Tx + \psi(x)^{2}N ,
\]
définie sans supposer $x', x''$, à valeurs dans $\mathrm{Sym}^{2}E$ mais en fait dans
$L^{\otimes 2}$ (« argument universel »), et dont la restriction à $L$ est
$u \mapsto u^{2}$. D'où une équivalence de catégories fibrées entre revêtements
quadratiques et triplets $(L, P_1, Q)$, $P_1$ un $L$-torseur et
$Q : P_1 \to L^{\otimes 2}$ une fonction polynôme de degré $2$ de partie homogène
$u \mapsto u^{2}$ ; le revêtement est le schéma des zéros de $Q$ dans le fibré
affine $P_1$.

La forme polaire de $N$ donne $\varphi_N : A \to E$, et
\[
  \boxed{\ Q(\varphi_N(x)) = -N(x)\,\delta\ }
\]
(vérifiable sur le cas décomposé, ou sur les formules de la page 52)\footnote{Le
signe est couvert page 46 ; page 52 il écrit $+N(x)\delta$, mais le calcul
qu'il y fait donne $-$, et la page 97 encadre $-N(x)\delta$.}. Sont équivalents :
$X'$ étale sur $X$ ; $\varphi_N$ isomorphisme ; $\delta$ inversible. Si $\delta = 0$,
$Q \circ \varphi_N = 0$.

\emph{Le cas où $2$ est inversible.} L'application $x \mapsto z = 2x - T$, $P_1 \to L$,
est un isomorphisme si $2$ est inversible, et
\[
  Q(T + z) = z^{2} - \delta, \qquad 4Q(x) = Q(2x) = Q(T + z).
\]
Une section $x$ de $X'$ donne donc une racine carrée $z$ de $\delta$, et
réciproquement quand $2$ est inversible : c'est la théorie de Kummer. L'involution
du revêtement est $x \mapsto \bar{x} = \psi(x)T - x$, et $\bar{z} = -z$\footnote{La
page 47 écrit « $\psi(T) = T$, $\psi(N) = N$ » ; la lettre est douteuse, et l'on lit
l'invariance de $T$ et $N$ par l'involution.}.

\emph{Ce qui reste quand $2$ ne l'est pas.} Sur $X_0 = V(2)$, $\psi(T) = 2 = 0$, donc
$T$ y est une section $T_0$ de $L_0$, et $T_0^{2} = \delta_0$. Mieux : tout relèvement
local $\tau$ de $T_0$ en une section de $L$ vérifie $\tau^{2} \equiv \delta$ modulo
$4L^{\otimes 2}$, indépendamment du relèvement, puisque $(\tau + 2a)^{2} \equiv \tau^{2}$
modulo $4$. Il l'écrit
\[
  \boxed{\ T_0^{2} \equiv \delta \pmod 4\ }.
\]
Si $x$ est une section de $X'$ et $z = 2x - T$, alors $z^{2} = \delta$ et $z_0 = T_0$ :
$z$ est une racine carrée de $\delta$ qui prolonge celle, $T_0$, déjà donnée sur
$X_0$.

\pagerange{48}{53}
\subsection*{Quand $2$ est régulier : les triplets $(L, \delta, T_0)$ (pages 48 à 53)}

Le foncteur $\mathfrak{X}'$ des racines carrées $z$ de $\delta$ telles que $z_0 = T_0$
est un faisceau fpqc, non représentable en général ; il est torseur, si $\delta$ est
inversible, sous $\mu'_2 = \{\lambda : \lambda^{2} = 1,\ \lambda_0 = 1\}$, et $f : X' \to
\mathfrak{X}'$, $x \mapsto 2x - T$, est compatible à $(\mathbf{Z}/2)_X \to \mu'_2$ :
$f(\bar{x}) = -f(x)$\footnote{La marge précise : non représentables s'il existe un
point où $2 \in \mathfrak{m}_x$ sans être nul dans $\mathcal{O}_{X,x}$.}.

\textbf{Théorème.} Supposons $2$ non diviseur de zéro dans $\mathcal{O}_X$. La donnée
d'un revêtement quadratique de $X$ équivaut à celle d'un triplet
\[
  (L, \delta, T_0), \qquad L \text{ inversible}, \quad \delta \in \Gamma L^{\otimes 2}, \quad
  T_0 \in \Gamma(X_0, L_0), \qquad T_0^{2} \equiv \delta \pmod 4 .
\]
Puisque $X - X_0$ est schématiquement dense, une section $x$ est connue par
$z = 2x - T$, et l'on reconstruit tout en divisant par $2$ ou $4$ là où c'est
permis. Pour $\tau$ un relèvement local de $T_0$ :
\[
  \Gamma E = \{(\lambda, \zeta) \mid \lambda \in \Gamma\mathcal{O},\ \zeta \in \Gamma(\lambda\tau + 2L)\},
  \qquad \lambda = \psi(x),\ \zeta = 2x - \psi(x)T ,
\]
\[
  \Gamma A = \{(b, \gamma) \in \Gamma\mathcal{O} \times \Gamma\check{L} \mid b_0 = \gamma_0 T_0\},
  \qquad b = \mathrm{Tr}(x),
\]
avec l'unité $(2, 0)$, la trace $T(b, \gamma) = b$, la norme et le produit
\[
  N(b, \gamma) = \frac{b^{2} - \gamma^{2}\delta}{4}, \qquad
  (b, \gamma)(b', \gamma') = \Bigl(\frac{bb' + \gamma\gamma'\delta}{2}, \frac{b\gamma' + \gamma b'}{2}\Bigr),
\]
le couplage $\langle (b, \gamma), (\lambda, \zeta)\rangle = \tfrac{1}{2}(\lambda b +
\gamma\zeta)$, l'application $\varphi_N(b, \gamma) = (b, -\gamma\delta)$, et la forme
\[
  Q(\lambda, \zeta) = \frac{\zeta^{2} - \lambda^{2}\delta}{4} .
\]
Chaque division est légitime : $b \equiv \gamma T$ et $\zeta \equiv \lambda T$ modulo $2$,
d'où $b^{2} - \gamma^{2}\delta \equiv \gamma^{2}(T^{2} - \delta) \equiv 0$ modulo $4$, et de
même pour $\zeta$ ; il le vérifie à chaque étape. La forme $Q$ est proportionnelle à
celle de la page 45, puisqu'elles ont les mêmes zéros sur $P_1$, et le facteur est
$1$ : en $y = T = (2, 0)$, les deux valent $-\delta$\footnote{La première lecture,
« $\xi \in \Gamma(\tfrac{\lambda\tau}{2} + L)$ », est remplacée par la description encadrée
en $\zeta = 2\xi$ ; la marge dit « ne dépend pas du choix de $L$ », pour « de
$\tau$ ». La lettre lue $\varphi$ page 50 est douteuse. Les signes de deux lignes de la
page 52 sont couverts ; on les donne d'après la formule encadrée.}.

Sur $\mathbf{Z}$, où $L = \mathcal{O}$, $X_0 = \operatorname{Spec}\mathbf{F}_2$ et
$T_0 \in \mathbf{F}_2$, la condition $T_0^{2} \equiv \delta$ modulo $4$ dit exactement
que $\delta \equiv 0$ ou $1$ modulo $4$, et que $T_0$ est déterminé par $\delta$ : le
théorème généralise la classification classique des anneaux quadratiques par leur
discriminant, avec la congruence de Stickelberger\footnote{Rapprochement de
l'édition.}.

\pagerange{53}{55}
\subsection*{Torseurs sous $\mu_2$ et sous $\mathbf{Z}/2$ (pages 53 à 55)}

L'homomorphisme $(\mathbf{Z}/2)_X \to \mu_{2,X}$ donne, si $2$ est régulier, le
diagramme

\begin{tikzcd}[column sep=small, nodes={font=\small}]
\mathrm{Tors}(\mathbf{Z}/2) \arrow[r, "\mathrm{Ass}"] \arrow[d, "\wr"] & \mathrm{Tors}(\mu_2) \arrow[r, "\mathrm{Oub}"] \arrow[d, "\wr"] & \mathrm{RevQuad}(X) \arrow[d, "\wr"] \\
\{(L, \delta, T_0),\ \delta \text{ inv.}\} \arrow[r, "\mathrm{Ass}'"] & \{(L, \theta),\ \theta \text{ inv.}\} \arrow[r, "\mathrm{Oub}'"] & \{(L, \delta, T_0)\}
\end{tikzcd}

avec $\mathrm{Ass}'(L, \delta, T_0) = (L, \delta)$ et $\mathrm{Oub}'(L, \theta) = (L, 4\theta, 0)$ :
le $\mu_2$-torseur $z^{2} = \theta$ vu comme revêtement quadratique, de discriminant
$4\theta$ et de trace nulle. $\mathrm{Oub}'$ ne commute pas aux produits, même si $2$
est inversible : $(LL', 16\theta\theta')$ n'est pas en général isomorphe à
$(LL', 4\theta\theta')$ --- sur $\mathbf{Z}$, avec $\theta = \theta' = 1$, il faudrait une
unité de carré $4$.

Les revêtements quadratiques forment au contraire une catégorie monoïdale
symétrique pour
\[
  (L, \delta, T_0) * (L', \delta', T'_0) = (L \otimes L', \delta\delta', T_0T'_0),
\]
qui prolonge le produit contracté des $(\mathbf{Z}/2)$-torseurs ; la condition
$(T_0T'_0)^{2} \equiv \delta\delta'$ modulo $4$ se vérifie aussitôt. C'est ce produit que
la suite du dossier construit sans l'hypothèse « $2$ régulier ».

\emph{Une catégorie mixte.} Soit $\mathrm{RevQuad}^{!}(X)$ la sous-catégorie pleine des
revêtements de $T_0 = 0$ : trace nulle sur $X_0$, norme additive sur $X_0$. Il
définit une catégorie monoïdale des couples $(n, \mathfrak{X})$, $n \in \mathbf{Z}/2$ une
parité, $\mathfrak{X} \in \mathrm{RevQuad}^{!}$ si $n = 1$, avec le produit usuel si
l'une des parités est paire, et, pour deux impaires,
\[
  (L, \delta, 0) *^{!} (L', \delta', 0) = (L \otimes L', \tfrac{1}{4}\delta\delta', 0) :
\]
$T_0 = 0$ force $\delta = 4\Delta$ localement, et $\tfrac{1}{4}\delta\delta' := 4\Delta\Delta'$ ne
dépend pas du choix de $\Delta$ à un $h$ près avec $4h = 0$\footnote{Si $2$ est
régulier, $h = 0$ et $\Delta$ est unique ; la latitude en $h$ prépare le cas
général. La marge de la page 54 dit que ce qui suit « doit être repris de façon
plus fine », en remplaçant les revêtements quadratiques par des
$\mu_2$-revêtements ; la fin en est illisible.}. Il reste à « lui donner un sens
indépendamment de la condition $2$ régulier, exercice catégorique universel » ;
le dossier ne le fait pas. C'est la même intention que celle des pages 1 et 2.

\pagerange{56}{57}
\subsection*{En caractéristique $2$ (pages 56 et 57)}

Si $2 = 0$ sur $X$, l'application $f_T : L \to L^{\otimes 2}$, $u \mapsto u^{2} + Tu$,
est additive (et semi-linéaire pour le Frobenius) ; elle est étale si et seulement
si $T$ est inversible, et s'identifie alors à l'application d'Artin--Schreier
$u \mapsto u^{2} + u$. Son noyau $G_T$ est un schéma en groupes fini localement
libre d'ordre $2$, étale si et seulement si $T$ est inversible\footnote{C'est un
des schémas en groupes d'ordre $2$ de la classification de Tate et Oort (1970) ;
le nom n'est pas sur la page.}. Si $x$ est une section de $X'$, $Q(x + u) = Q(x) +
u^{2} + Tu$ ; donc $X'$ est un torseur sous $G_T$.

\textbf{Théorème.} Sur les schémas de caractéristique $2$ (ou les topos localement
annelés de caractéristique $2$), les revêtements quadratiques équivalent aux
données $(L, P_1, T, \alpha)$ : $L$ inversible, $P_1$ un $L$-torseur, $T \in \Gamma L$,
$\alpha$ une section du $L^{\otimes 2}$-torseur $f_{T*}P_1$ déduit de $P_1$ par $f_T$.
Le revêtement est la fibre en $\alpha$ de l'application canonique
$P_1 \to f_{T*}P_1$, torseur sous $G_T$ : « Plus de doute, c'est lui ! »\footnote{En
marge : l'application $P_1 \to P_1^{(2)}$, $x \mapsto Tx + N$, compatible à $u \mapsto
Tu$, et l'équation $Q(x) = 0$ s'écrit $x^{2} = Tx + N$.}

Le produit contracté s'y lit sans peine : $(u, u') \mapsto uT' + Tu'$ est un
homomorphisme $G_T \times G_{T'} \to G_{TT'}$, puisque $f_{TT'}(uT') = T'^{2}f_T(u)$, et
l'extension du groupe structural donne $\mathfrak{X} * \mathfrak{X}'$, torseur sous
$G_{TT'}$.

\pagerange{58}{61}
\subsection*{Quatre catégories équivalentes (pages 58 à 61)}

Sur les schémas, ou sur les topos localement annelés, sont canoniquement
équivalentes :
\begin{enumerate}
\item[(1)] les revêtements quadratiques $X'/X$ ; (1') les $\mathcal{O}$-algèbres
commutatives localement libres de rang $2$ ;
\item[(2)] les triplets $(A, N, e)$, $(A, N)$ un module quadratique localement libre de
rang $2$, $e$ une section avec $N(e) = 1$ ;
\item[(3)] les triplets $(L, P, Q)$, $Q : P \to L^{\otimes 2}$ polynomiale de partie
homogène $u \mapsto u^{2}$, c'est-à-dire $Q(x + u) = Q(x) + u^{2} + \ell_x(u)$ avec
$\ell_x$ affine en $x$ ;
\item[(3')] les couples formés d'une extension $0 \to L \to E \xrightarrow{\psi}
\mathcal{O} \to 0$, $L$ inversible, et d'une forme $Q : E \to L^{\otimes 2}$ de
restriction $u \mapsto u^{2}$ à $L$.
\end{enumerate}
La fonction $Q(x) - x^{2}$ sur $P$ est affine (sa partie quadratique est nulle).
Dans un scindage $e \in \Gamma P_1$, $Q(e + u) = u^{2} + bu + c$, et
\[
  Q(x) = x^{2} - xT + N, \qquad \boxed{T = 2e - b, \quad N = e^{2} + c - be}, \qquad \psi(T) = 2,\ N(e) = 1,
\]
indépendants du scindage. Le couple $(T, N)$ est un point d'un torseur sous
$L^{\otimes 2} \otimes \check{E} \simeq L \otimes E$ (pour $E$ de rang $2$,
$E \simeq \check{E} \otimes \det E$), l'espace des applications affines $P_1 \to
L^{\otimes 2}$, celle définie par $\xi \in L \otimes E$ étant $x \mapsto x \wedge \xi$.
Inversement, $(T, N)$ provient d'une $Q$ si et seulement si $\psi(T) = 2$, $N(e) = 1$, et
$T(x) = \varphi_N(x, e)$ : « $N$ détermine $T$ ». La condition $4N - T^{2} \in
L^{\otimes 2}$ ne suffit pas ; en coordonnées, $T = 2x_0 + t$ et $N = x_0^{2} - bx_0 + c$
exigent $t = -b$, et alors
\[
  Q(x_0 + u) = u^{2} + bu + c, \qquad \delta = T^{2} - 4N = b^{2} - 4c .
\]

\pagerange{62}{65}
\subsection*{Le produit contracté, deviné (pages 62 à 65)}

On veut un produit $X_1 * X_2$ avec $L(X_1 * X_2) = L_1 \otimes L_2$, $\delta = \delta_1\delta_2$,
et une application $(x_1, x_2) \mapsto x_1 * x_2$, $P_{X_1} \times P_{X_2} \to P_{X_1 *
X_2}$, qui envoie $X_1 \times X_2$ dans $X_1 * X_2$. Dans des scindages, il pose
$b = b_1b_2$ (« sauf erreur ») et en tire $c$ par $4c = b^{2} - \delta$ :
\[
  \boxed{\ c = c_1\delta_2 + c_2\delta_1 + 4c_1c_2 = b_2^{2}c_1 + b_1^{2}c_2 - 4c_1c_2\ }
\]
--- expression qui, sous sa seconde forme, n'a plus de dénominateur\footnote{L'encadré
de la page 63 écrit « $4c_1c_1$ » ; le calcul et le membre de droite donnent
$4c_1c_2$. Le tableau de la page 62 mêle les indices ; on lit $b_i \in \Gamma L_i$,
$c_i \in \Gamma L_i^{\otimes 2}$.}. Puis : « Je devine qu'on a »
\[
  (x_1 + u_1) * (x_2 + u_2) = x_1 * x_2 + (u_1b_2 + u_2b_1 + 2u_1u_2).
\]
Il le vérifie par l'identité, valable pour tous $u_i \in \Gamma L_i$, où
$Q_i(u_i) = u_i^{2} + b_iu_i + c_i$ et $u = u_1b_2 + u_2b_1 + 2u_1u_2$ :
\[
  u^{2} + bu + c = b_2^{2}Q_1(u_1) + b_1^{2}Q_2(u_2) + 4\bigl[Q_1(u_1)Q_2(u_2) - c_1Q_2(u_2) - c_2Q_1(u_1)\bigr],
\]
qui s'annule quand $Q_1(u_1) = Q_2(u_2) = 0$\footnote{On a refait le développement :
il est exact. La page porte des essais biffés ; l'accolade « $-4c_1c_2$ » sous
$4(c_1c_2 - c_1c_2 - c_1c_2)$ est la sienne.}. Il en retient, sous forme
intrinsèque,
\[
  \boxed{\ Q(x_1 * x_2) = Q_1(x_1)\,\delta_2 + Q_2(x_2)\,\delta_1 + 4\,Q_1(x_1)Q_2(x_2)\ } .
\]
L'identité ci-dessus montre que cette formule, si elle vaut pour un couple de
points de base, vaut pour tous : son second membre se réécrit
$Q_1\delta_2 + Q_2\delta_1 + 4Q_1Q_2$ avec $Q_i = Q_i(x_i + u_i)$.

\pagerange{65}{79}
\subsection*{Digression affine : les torseurs $\chi$-trivialisés (pages 65 à 79)}

Soit $\chi \in \Gamma\mathcal{O}$ fixé. Il considère les extensions
$0 \to L \to E \xrightarrow{\psi} \mathcal{O} \to 0$ (c'est-à-dire les $L$-torseurs
$P = \psi^{-1}(1)$) munies d'une \emph{$\chi$-trivialisation}, donnée sous l'une des
trois formes équivalentes :
\begin{itemize}
\item une section $T$ de $E$ avec $\psi(T) = \chi$ (une trivialisation de
$P \wedge^{L}(L \xrightarrow{\chi} L)$) ;
\item une « $\chi$-section » $\mathcal{O} \to E$, $\lambda \mapsto \lambda T$ ;
\item une « $\chi$-rétraction » $\pi : E \to L$, $\pi|_L = \chi\,\mathrm{id}_L$.
\end{itemize}
Elles se correspondent par
\[
  \boxed{\ \pi(x) = \chi x - \psi(x)\,T\ } ;
\]
en effet $\psi(\chi x - \psi(x)T) = 0$, et réciproquement $\chi x - \pi(x)$, nul sur $L$,
se factorise par $E/L = \mathcal{O}$\footnote{À partir de la page 66, le module noté
$M$ page 65 devient $L$ ; plusieurs $M$ sont surchargés en $L$. Page 67, « section
locale $x$ de $U$ » se lit « de $E$ ».}. Pour $\chi = 2$ et $E = \check{A}$, $T$ est la
trace, et $b = \pi(x_0) = 2x_0 - T$ est le $b$ des scindages.

\emph{Le produit.} Pour deux telles extensions, soit $E_1 \otimes_0 E_2 = L_1 \otimes E_2 +
E_1 \otimes L_2 \subset E_1 \otimes E_2$, noyau de $\psi_1 \otimes \psi_2$, et
$\mu : E_1 \otimes_0 E_2 \to L_1 \otimes L_2$ défini par $\mathrm{id} \otimes \pi_2$ sur
$L_1 \otimes E_2$ et $\pi_1 \otimes \mathrm{id}$ sur $E_1 \otimes L_2$ --- qui coïncident,
égaux à $\chi$, sur $L_1 \otimes L_2$. Le produit $E_1 * E_2$ est la somme amalgamée :

\begin{tikzcd}
0 \arrow[r] & E_1 \otimes_0 E_2 \arrow[r] \arrow[d, "\mu"] & E_1 \otimes E_2 \arrow[r, "\psi_1 \otimes \psi_2"] \arrow[d] & \mathcal{O} \arrow[r] \arrow[d, no head, "=" description] & 0 \\
0 \arrow[r] & L_1 \otimes L_2 \arrow[r] & E_1 * E_2 \arrow[r, "\psi"] & \mathcal{O} \arrow[r] & 0
\end{tikzcd}

et sa $\chi$-rétraction est
\[
  \boxed{\ \pi(x_1 * x_2) = \pi_1(x_1) \otimes \pi_2(x_2)\ } ,
\]
bien définie sur la somme amalgamée parce que $\pi_1 \otimes \pi_2$ et
$\chi\,\mathrm{id}_{L_1 \otimes L_2}$ coïncident sur $E_1 \otimes_0 E_2$ à travers $\mu$
(page 70 : « On y gagne ! »). Aucune division n'est nécessaire. Si $\chi$ est
régulier sur $L_1 \otimes L_2$, la section correspondante est $T = \tfrac{1}{\chi}\,T_1 *
T_2$ ; en général on a $\chi T = T_1 * T_2$, et, pour $x_i \in \Gamma P_i$,
\[
  T = \chi\,(x_1 * x_2) - b_1 \otimes b_2, \qquad b_i = \pi_i(x_i).
\]\footnote{Les pages 71 et 72
  hésitent sur le signe de $b$ (« attention au signe ! ») et le fixent par
  $b_{T, x} = \chi x - T = \pi(x)$, convention suivie ici. La note de la page 79
  écrit $T = 2\chi x - b$ ; c'est $\chi x_0 - b$.}

\emph{Formulaire} (pages 72 à 74). L'application $(x_1, x_2) \mapsto x_1 * x_2$ est
bilinéaire, $\psi(x_1 * x_2) = \psi_1(x_1)\psi_2(x_2)$, donc elle envoie $P_1 \times P_2$
dans $P$, et
\[
  u_1 * x_2 = u_1 \otimes \pi_2(x_2), \qquad x_1 * u_2 = \pi_1(x_1) \otimes u_2, \qquad
  u_1 * u_2 = \chi\, u_1 \otimes u_2,
\]
\[
  (x_1 + u_1) * (x_2 + u_2) = x_1 * x_2 + u_1 \otimes \pi_2(x_2) + \pi_1(x_1) \otimes u_2 + \chi\, u_1 \otimes u_2 .
\]
Pour $\chi = 2$, c'est la formule devinée page 63.

\emph{Propriété universelle} (pages 75 et 76). Une application « bi-affine »
compatible aux $\chi$-trivialisations, de $P_1 \times P_2$ dans un torseur $\Pi$ sous
$M$, est un couple $(f, \alpha)$, $\alpha : L_1 \times L_2 \to M$ bilinéaire, avec
\[
  f(x_1 + u_1, x_2) = f(x_1, x_2) + \alpha(u_1, \pi_2(x_2)), \qquad
  f(x_1, x_2 + u_2) = f(x_1, x_2) + \alpha(\pi_1(x_1), u_2) ;
\]\footnote{La page écrit
  $\alpha(x_1, x_2)$ pour premier terme, où l'on attend $f(x_1, x_2)$ ; chaque autre
  $\alpha$ est écrit sur une majuscule biffée.}
la donnée de $\alpha$ n'est pas superflue, car $(E_1 \otimes L_2) \oplus (L_1 \otimes E_2)
\to L_1 \otimes L_2$ n'est pas surjectif en général. $(P_1 * P_2, *)$ est universel
pour ces applications, et le reste si l'on exige aussi $\pi_\Pi(f(x_1, x_2)) =
\pi_1(x_1) \otimes \pi_2(x_2)$. Un morphisme de torseurs $\chi$-trivialisés est un
morphisme affine $f$, de partie linéaire $f_0$, avec $\pi'(f(x)) = f_0(\pi(x))$ ;
en termes d'extensions, $f(T) = T'$.

\emph{Le cas scindé} (pages 77 à 79). Un scindage est une section $x_0$ de $P$ ; la
$\chi$-trivialisation équivaut alors à $b = \pi(x_0) \in \Gamma L$, avec
$T = (-b, \chi)$ et $\pi(u, \lambda) = \lambda b + \chi u$ dans $E \simeq L \oplus
\mathcal{O}$. Pour des scindages $x_1, x_2$, $P_1 * P_2$ est scindé par $x_1 * x_2$ et
\[
  b = b_1 \otimes b_2
\]
--- « Il n'y a pas plus simple ! ». Les contraintes de commutativité et
d'associativité sont caractérisées par $x_1 * x_2 \mapsto x_2 * x_1$, $(x_1 * x_2) * x_3
\mapsto x_1 * (x_2 * x_3)$ et les isomorphismes correspondants des $L_i$. L'unité est
$E_0 = \mathcal{O} \times \mathcal{O}$, $x_0 = (0, 1)$, $\pi(u, \lambda) = \lambda + \chi u$,
$T = (-1, \chi)$ : un scindage est \emph{unitaire} si $b$ est une base de $L$, et
$P_0 * P_1 \simeq P_1$ par $x_1 \mapsto x_0 * x_1$, $u_1 \mapsto b_0 \otimes u_1$. Un objet
est inversible pour $*$ si et seulement si $L$ est inversible et, sur $X_0 =
V(\chi)$, la section $T_0 = -b_0$ de $L_0$ est une base\footnote{Cette
caractérisation est dans une longue note marginale, oblique, en partie illisible ;
on donne ce qui s'en lit, qui est juste : si $b_0$ engendre $L_0$, $b$ engendre $L$
près de $X_0$, et hors de $X_0$, où $\chi$ est inversible, on rend $b$ unitaire par
$b \mapsto b + \chi u$. La nécessité : $b \otimes b'$ unitaire force $b_0$ à être une
base sur $X_0$, où $b'$ ne change pas.}. Pour $\chi = 2$ : l'extension
$2$-trivialisée d'un revêtement quadratique est inversible si et seulement si le
revêtement est étale en tout point de caractéristique $2$.

\pagerange{79}{87}
\subsection*{Le cas $\chi = 2$ : involution et coinvariants (pages 79 à 87)}

« Les revêtements quadratiques : à partir de maintenant $\chi = 2$. » L'involution
$\sigma_E(x) = \bar{x} = \psi(x)T - x$ respecte $\psi$, induit $-\mathrm{id}$ sur $L$ et
$\mathrm{id}$ sur $\mathcal{O}$. Si $2$ est inversible, $x_0 = \tfrac{1}{2}T$ scinde
canoniquement, $b = 0$, et le produit $*$ se réduit au produit tensoriel des $L_i$ :
« La théorie est tt.\ triviale pour $2$ inversible ! […] Le cas intéressant est
celui où $2$ n'est pas inversible. » On a
\[
  \overline{x_1 * x_2} = \bar{x}_1 * x_2 = x_1 * \bar{x}_2, \qquad \bar{x}_1 * \bar{x}_2 = x_1 * x_2,
\]
par exemple $\bar{x}_1 * x_2 = x_1 * x_2 - b_1 \otimes b_2 = T - x_1 * x_2$.

L'application $E_1 \otimes E_2 \to E_1 * E_2$ se factorise donc par les coinvariants de
$\sigma_1 \otimes \sigma_2$. Quand est-ce un isomorphisme ?

\textbf{Proposition} (page 85). Supposons $L_1$, $L_2$ inversibles. Pour que $E_1 \otimes E_2
\to E_1 * E_2$ induise un isomorphisme
\[
  (E_1 \otimes E_2)/\operatorname{Im}(\sigma_1 \otimes \sigma_2 - \mathrm{id}) \xrightarrow{\ \sim\ } E_1 * E_2,
\]
il faut et il suffit qu'en tout point de $X_0 = V(2)$, $T_{1,0}$ ou $T_{2,0}$ engendre
$L_{1,0}$, resp.\ $L_{2,0}$ --- c'est-à-dire que $X_1$ ou $X_2$ y soit étale. Avec
des bases et des scindages, la condition s'écrit $2\mathcal{O} + b_1\mathcal{O} +
b_2\mathcal{O} = \mathcal{O}$.

La condition est celle de la surjectivité : l'image de $\pi_i$ est
$\mathcal{O}b_i + 2L_i$, et le conoyau de $E_1 \otimes E_2 \to E_1 * E_2$ est
$\operatorname{Coker}\pi_1 \otimes \operatorname{Coker}\pi_2$. L'égalité du noyau et de
l'image de $\sigma_1 \otimes \sigma_2 - \mathrm{id}$ s'en déduit. La page le vérifie au-dessus
de $X_0$ par un décompte de rangs ; on peut aussi le faire directement.
Localement, $L_i = \mathcal{O}e_i$, $b_i = \beta_i e_i$, et près d'un point où $\beta_1$ est
inversible : dans la base $(e_1 \otimes e_2, e_1 \otimes x_2, x_1 \otimes e_2, x_1 \otimes x_2)$,
le noyau est $\{(a, a', a'', 0) : 2a + \beta_2a' + \beta_1a'' = 0\}$, de base
$(1, 0, -2/\beta_1, 0)$ et $(0, 1, -\beta_2/\beta_1, 0)$, et l'image de $\sigma_1 \otimes
\sigma_2 - \mathrm{id}$ contient $(\beta_1, 0, -2, 0)$ et
$(\beta_1\beta_2, -\beta_1, -\beta_2, 0)$, qui en sont, à des unités près, une
base\footnote{Ce calcul est de l'édition ; il évite la réduction au cas où
$2 = 0$ sur $X$, que la page annonce sans la justifier et qui demanderait un
argument sur une base non réduite. Les générateurs sont ceux de la page 84. La
page 85 appelle $n$ le rang de $L_1 \otimes L_2$, qui vaut $1$. Sans la condition,
par exemple si $T_1 = T_2 = 0$, $\sigma_i = \mathrm{id}$ et l'image est nulle alors que
le noyau est de rang $3$ (le crochet de la page 84).}.

\textbf{Corollaire} (page 86). Sous la même condition, et pour tout Module $M$,
$\underline{\mathrm{Hom}}(E_1 * E_2, M) \simeq \underline{\mathrm{Hom}}(E_1 \otimes
E_2, M)^{\sigma_1 \otimes \sigma_2}$ ; en particulier, avec $A_i = \check{E}_i$ et
$A = (E_1 * E_2)^{\vee}$,
\[
  A_1 * A_2 \xrightarrow{\ \sim\ } (A_1 \otimes A_2)^{\sigma_1 \otimes \sigma_2},
\]
et réciproquement si cela reste vrai après tout changement de base. C'est la
définition classique du produit des algèbres quadratiques étales ; elle reste
valable dès qu'en chaque point de caractéristique $2$ l'un des facteurs est
étale\footnote{Pour le produit des algèbres quadratiques étales, voir M.-A. Knus,
\emph{Quadratic and Hermitian Forms over Rings} (1991), chap.~III ; pour des algèbres
quadratiques quelconques, la question a été reprise par d'autres moyens dans les
travaux sur l'« algèbre discriminante » (Deligne, Rost, Loos, Biesel--Gioia) ;
on ne prétend pas qu'ils retrouvent exactement cette construction. Références de
l'édition.}.

\emph{Au-dessus de $V(\chi)$} (pages 86 à 88). Si $\chi = 0$ sur $X$, une
$\chi$-trivialisation n'est plus qu'une section $T$ de $L$, et $\pi(x) = -T$ est
constant. Alors $u_1 * u_2 = 0$, $E_1 \otimes E_2 \to E_1 * E_2$ se factorise par
$E_1 \otimes E_2/L_1 \otimes L_2 \simeq E_1 \times_{\mathcal{O}} E_2$, l'extension du torseur
produit $P_1 \times P_2$, et $E_1 * E_2$ s'en déduit par $(u_1, u_2) \mapsto u_1 \otimes b_2
+ b_1 \otimes u_2$ : $P_1 * P_2$ est poussé de $P_1 \times P_2$ par cet homomorphisme,
comme $G_T \times G_{T'} \to G_{TT'}$ page 57\footnote{La page note le premier membre
d'un $\otimes$ surmonté de deux demi-flèches.}.

\pagerange{88}{92}
\subsection*{La forme quadratique du produit (pages 88 à 92)}

Soient maintenant sur $E_1, E_2$ des formes $Q_i : E_i \to L_i^{\otimes 2}$, de restriction
$u \mapsto u^{2}$ et de trace $T_i$ (des revêtements quadratiques $X_1, X_2$), de
discriminants $\delta_i = b_i^{2} - 4c_i$, avec, dans des scindages,
$Q_i(\lambda x_i + u_i) = u_i^{2} + \lambda b_iu_i + \lambda^{2}c_i$.

\textbf{Théorème} (page 89). Il existe sur $E = E_1 * E_2$ une unique forme
quadratique $Q$ à valeurs dans $L^{\otimes 2}$, $L = L_1 \otimes L_2$, admissible ---
de restriction $u \mapsto u^{2}$ à $L$ et de trace la section $T$ de la
$2$-trivialisation de $E$ --- telle que, pour $x_i \in \Gamma P_i$,
\[
  Q(x_1 * x_2) = Q_1(x_1)\,\delta_2 + Q_2(x_2)\,\delta_1 + 4\,Q_1(x_1)Q_2(x_2),
\]
et, pour $x_i \in \Gamma E_i$ quelconques, avec des facteurs $\psi_2(x_2)^{2}$,
$\psi_1(x_1)^{2}$ devant les deux premiers termes. Dans les scindages $x_1, x_2$,
$x_1 * x_2$ :
\[
  \boxed{\ b = b_1b_2, \qquad c = c_1b_2^{2} + c_2b_1^{2} - 4c_1c_2\ } .
\]
\textbf{Corollaire.} $\delta = \delta_1\delta_2$ : $b_1^{2}b_2^{2} - 4(c_1b_2^{2} + c_2b_1^{2} -
4c_1c_2) = (b_1^{2} - 4c_1)(b_2^{2} - 4c_2)$.

Les deux formules simples sont $b = b_1b_2$ et $\delta = \delta_1\delta_2$ ; celle de $c$
s'en déduit « heuristiquement » en divisant $4c = b^{2} - \delta$ par $4$, et se
démontre localement par l'identité de la page 64, « un petit calcul délicat, que
j'ai pratiquement déjà fait plus haut »\footnote{La page définit « admissible » par
la seule condition sur $L$ ; l'unicité demande aussi que la trace de $Q$ soit le
$T$ de $E$, ce qu'il utilise en posant $b = b_1b_2$. La note marginale de la page
88, qui récapitule la notion de forme sur une extension $2$-trivialisée, n'est lisible
qu'en débris ; le texte courant de cette page l'est en partie seulement.}. En
termes de la page 46 : $b = \pi(x) = 2x - T$ est la « racine carrée » $z$ de
$\delta$ attachée au point $x$ ; la formule $b = b_1 \otimes b_2$ dit que le produit des
revêtements multiplie les racines carrées, $z = z_1z_2$, et le reste est l'arithmétique
qui rend cela vrai sans diviser par $2$\footnote{Cette lecture est de l'édition.}.

\emph{Un programme} (pages 91 et 92). Définir $X_1 * X_2$ et $X_1 \times X_2 \to X_1 * X_2$ ;
les contraintes de commutativité, d'associativité, d'unité ; décrire $X_1 * X_2$
quand $X_1$ est étale comme $X_2$ « tordu » par le $(\mathbf{Z}/2)$-torseur $X_1$ ;
étudier quand $X_1 * X_2$ est étale ; et l'énoncé « $X_1 \times X_2 \to X_1 * X_2$ est
fid.\ plat », dont la suite est illisible. Tel quel, ce dernier est faux : en
caractéristique $2$, pour $X_1 = X_2 = \operatorname{Spec}\mathcal{O}[u]/(u^{2})$, on a
$b_1 = b_2 = 0$, l'application est $u \mapsto 2u_1u_2 = 0$, et
$\mathcal{O}[u_1, u_2]/(u_1^{2}, u_2^{2})$ n'est pas plat sur $\mathcal{O}[u]/(u^{2})$ ;
il est vrai sous la condition de la proposition de la page 85\footnote{Le
contre-exemple est de l'édition. La page 92, barrée, commence un calcul en
coordonnées et évoque des revêtements « où $X_1 * X_2$ est non étale ».}.

\pagerange{93}{97}
\subsection*{Dualité entre l'algèbre et l'extension (pages 93 à 97)}

Un scindage $x_0$ donne $E \simeq L \oplus \mathcal{O}x_0$ et $A \simeq \mathcal{O} \oplus
L^{\vee}$, où $L^{\vee} = \operatorname{Ker}(x_0)$. Pour $u \in L^{\vee}$, $\langle u, b\rangle
= -\mathrm{Tr}(u)$, $N(u) = \langle u^{\otimes 2}, c\rangle$, d'où l'équation minimale
$u^{2} + \langle u, b\rangle u + \langle u^{\otimes 2}, c\rangle = 0$\footnote{La page
écrit $u^{2} + \langle u, b\rangle + \langle u^{\otimes 2}, c\rangle = 0$, sans le facteur
$u$ du terme du milieu.}, et $N = x_0^{2} - x_0b + c$. Une section $x = x_0 + v$ de
$P_1$ est un homomorphisme d'algèbres $A \to \mathcal{O}$ si et seulement si
$v^{2} + bv + c = 0$, c'est-à-dire $Q(v) = 0$.

\emph{L'isomorphisme $g$.} Pour $E$ de rang $2$, $x \mapsto (y \mapsto x \wedge y)$ donne
$E \simeq \underline{\mathrm{Hom}}(E, \det E)$, et $\alpha(u) = x_0 \wedge u$, $L \simeq \det E$,
ne dépend pas de $x_0$. D'où $g : E \xrightarrow{\sim} A \otimes L$, $g(x) = f_x$,
\[
  f_x(y) = \psi(x)\,y - \psi(y)\,x ,
\]
vérifié en évaluant sur $x = u$, $y = x_0$. Et $g$ est un isomorphisme d'extensions

\begin{tikzcd}
0 \arrow[r] & L \arrow[r] \arrow[d, "-\mathrm{id}_L"'] & E \arrow[r, "\psi"] \arrow[d, "g"] & \mathcal{O} \arrow[r] \arrow[d, no head, "=" description] & 0 \\
0 \arrow[r] & L \arrow[r] & A \otimes L \arrow[r] & \mathcal{O} \arrow[r] & 0
\end{tikzcd}

--- « attention au signe » : sur $L$, $g$ induit $-\mathrm{id}_L$. Par $g$,
$\underline{\mathrm{Quad}}(E, L^{\otimes 2}) \simeq \underline{\mathrm{Quad}}(A, \mathcal{O})$,
et \emph{$Q$ et $N$ se correspondent} ; le signe de $-\mathrm{id}_L$ change celui du
terme en $b$ : $Q(u + \lambda x_0) = u^{2} + \lambda\langle u, b\rangle + \lambda^{2}c$ et
$N(v + \lambda) = \lambda^{2} - \lambda\langle v, b\rangle + \langle v^{2}, c\rangle$. La
formule $Q(\varphi_N(x)) = -N(x)\,\delta$ est une identité entre formes sur le module
générique de rang $2$ muni d'une forme quadratique\footnote{Le signe de
$-\lambda\langle v, b\rangle$ est repris sur une rature, d'une lecture conjecturale ; c'est
celui que donne la page 93.}.

\pagerange{98}{106}
\subsection*{Trace et norme de $A_1 \otimes A_2$ sur $A_1 * A_2$ (pages 98 à 106)}

Transposé, $E_1 \otimes E_2 \to E_1 * E_2$ donne $i : A = A_1 * A_2 \to A_1 \otimes A_2$, et
la somme amalgamée devient, côté algèbres :
\[
  A = (A_1 \otimes A_2) \amalg_{A_1 \oplus_{\mathcal{O}} A_2} \mathcal{O}, \qquad
  \tau = (\mathrm{Tr}_1, \mathrm{Tr}_2) : A_1 \oplus_{\mathcal{O}} A_2 \to \mathcal{O},
\]
extension de $L_1^{\vee} \otimes L_2^{\vee}$ par $\mathcal{O}$ ($\mathrm{Tr}_1(1) = \mathrm{Tr}_2(1)
= 2$ rend $\tau$ bien défini). En bases $(1, U_i)$, $U_i^{2} + b_iU_i + c_i = 0$ :
\[
  A \simeq \mathcal{O}[\mathcal{U}]/(\mathcal{U}^{2} + b\,\mathcal{U} + c), \qquad
  i(\mathcal{U}) = 2U_1U_2 + b_1U_2 + b_2U_1,
\]
avec $b, c$ comme au théorème de la page 89 ; que $i(\mathcal{U})$ vérifie l'équation
est l'identité de la page 64, dans l'anneau $A_1 \otimes A_2$.

\textbf{Proposition} (pages 104 et 105). 1) L'homomorphisme canonique
$\mathrm{Tr}_{A_1 \otimes A_2/A} : A_1 \otimes A_2 \to A$ (le transposé de $i$), qu'il note
aussi $x_1 \otimes x_2 \mapsto x_1 * x_2$, vérifie
\[
  i\bigl(\mathrm{Tr}_{A_1 \otimes A_2/A}(\xi)\bigr) = \xi + \bar{\xi}, \qquad \bar{\xi} = (\sigma_1 \otimes \sigma_2)(\xi);
\]
sur la base : $1 \mapsto 2$, $U_i \mapsto -b_i$, $x_1 * 1 = \mathrm{Tr}_1(x_1)$, et
$U_1 * U_2 = \mathcal{U} + b_1b_2$, puisque $U_1U_2 + \bar{U}_1\bar{U}_2 = i(\mathcal{U}) + b_1b_2$.
2) Il existe une unique application quadratique $N_{A_1 \otimes A_2/A} : A_1 \otimes A_2 \to
A$ avec $N(x_1 \otimes x_2) = N_1(x_1)N_2(x_2)$ et $\varphi_N(\xi, \eta) =
\mathrm{Tr}_{A_1 \otimes A_2/A}(\xi\bar{\eta})$ ; elle vérifie $i(N(\xi)) = \xi\bar{\xi}$,
$N(\bar{\xi}) = N(\xi)$, $N((\sigma_1 \otimes \mathrm{id})\xi) = N((\mathrm{id} \otimes
\sigma_2)\xi) = \overline{N(\xi)}$, $N(i(\eta)) = \eta^{2}$, et la transitivité
\[
  N_{A/\mathcal{O}} \circ N_{A_1 \otimes A_2/A} = N_{A_1 \otimes A_2/\mathcal{O}}, \qquad
  \mathrm{Tr}_{A/\mathcal{O}} \circ \mathrm{Tr}_{A_1 \otimes A_2/A} = \mathrm{Tr}_{A_1 \otimes A_2/\mathcal{O}}.
\]\footnote{Page
  105, le second terme de la deuxième ligne s'écrit avec $\mathrm{id}_{A_2}$ ; on lit
  $\mathrm{id}_{A_1}$. L'existence de $N$ « universellement » est affirmée, non
  démontrée : on la définit sur la base et l'on vérifie les propriétés.}
Pour la transitivité de la norme, avec $\sigma_1 = \sigma_{A_1} \otimes \mathrm{id}$,
$\sigma_2 = \mathrm{id} \otimes \sigma_{A_2}$, $\sigma_3 = \sigma_1\sigma_2$ et $\xi_j =
\sigma_j(\xi)$, on a $N_{A_1 \otimes A_2/\mathcal{O}}(\xi) = \xi\xi_1\xi_2\xi_3$ (norme de
$A_1 \otimes A_2$ sur $A_1$, puis de $A_1$ sur $\mathcal{O}$), et ce produit est
$N\bar{N}$ pour $N = N_{A_1 \otimes A_2/A}(\xi)$. Si $A \to A_1 \otimes A_2$ n'est pas
injectif, l'argument est générique\footnote{La page note $\xi^{\natural} = \xi + \bar{\xi}$
(un signe en crochet, lu comme un bécarre), et remarque (page 101) que si, en tout
point de $X_0$, l'une des deux structures est étale, $A \to A_1 \otimes A_2$ est
universellement injectif d'image $(A_1 \otimes A_2)^{\sigma_1 \otimes \sigma_2}$ : c'est le
corollaire de la page 86. Le début de la page 101 est en partie illisible.}.

\pagerange{106}{109}
\subsection*{Algèbres quasi-quadratiques (pages 106 à 109)}

« Quand on axiomatise la situation » de $\mathrm{Tr}_{A_1 \otimes A_2/A}$ et
$N_{A_1 \otimes A_2/A}$, on aboutit à une $k$-algèbre $A$ « qui se comporte comme si
c'était une algèbre quadratique sur $k$ », sans être supposée localement libre de
rang $2$, ni même avoir $\alpha : k \to A$ injectif. Une \emph{algèbre
quasi-quadratique} sur $k$ est la donnée de $A$, $\alpha$, d'une application
$k$-linéaire $t : A \to k$ et d'une application $n : A \to k$, avec, en posant
$\tau = \alpha \circ t$, $\nu = \alpha \circ n$, $\bar{x} = \tau(x) - x$ :
\begin{enumerate}
\item[(1)] $t(1_A) = 2$ ;
\item[(2)] $x \mapsto \bar{x}$ est un automorphisme d'algèbre, c'est-à-dire
$2xy = \tau(xy) - \tau(x)\tau(y) + \tau(y)x + \tau(x)y$ ;
\item[(3)] $n$ est quadratique, de forme polaire $\varphi_n(x, y) = t(x\bar{y})$ ;
\item[(4)] $n(1) = 1$, ce qui, avec (3), redonne $t(1) = 2$ ;
\item[(5)] $\nu(x) = x\bar{x}$, ce qui implique (3) et (4) si $\alpha$ est injectif ;
\item[(6)] facultatif : si $k$ est fini localement libre de rang $n$ sur $k_0$ et $A$
de rang $2n$, $\mathrm{Tr}_{k/k_0} \circ t = \mathrm{Tr}_{A/k_0}$ et $N_{k/k_0} \circ n =
N_{A/k_0}$.
\end{enumerate}
L'involution est involutive par (1) ; $n(\bar{x}) = n(x)$ par (3). La page 107 inscrit
aussi la multiplicativité $n(xy) = n(x)n(y)$, que le récapitulatif de la page 109
omet ; avec (5) et $\alpha$ injectif, elle en découle\footnote{Page 108, au point
4°, la page dit « $k_0 \to A$ injectif mais $k \to A$ injectif » ; le sens est
« mais non nécessairement $k \to A$ ». Le facteur $2$ de (2), remarque-t-il, n'est
gênant que si $2$ n'est pas régulier : la structure d'algèbre est alors connue par
$1_A$, $\alpha$, $t$ et $n$. La notion est voisine de celle d'algèbre munie d'une
involution standard (Voight) ; rapprochement de l'édition.}.

\pagerange{110}{121}
\subsection*{La situation quasi-biquadratique (pages 110 à 121)}

Données : $k$ ; des $k$-algèbres quasi-quadratiques $A_1, A_2, A_3$ ; une
$k$-algèbre $A$ avec des homomorphismes $\alpha_i : A_i \to A$, non nécessairement
injectifs ; et sur $A$ des structures quasi-quadratiques $(T_i, N_i)$ au-dessus de
chaque $A_i$, d'involutions $\sigma_i(x) = T_i(x) - x$. Compatibilités :
\begin{enumerate}
\item[(1)] la situation est, pour les traces, celle d'un changement de base :
$T_2$ est $t_1 \otimes \mathrm{id}_{A_2}$ sur l'image de $A_1 \otimes_k A_2$, ce qui exige
que $t_1 \otimes \mathrm{id}$ s'annule sur le noyau de $A_1 \otimes A_2 \to A$ ; donc
$\sigma_2$ et $\sigma_3$ induisent sur $A_1$ son involution ;
\item[(2)] de même pour les normes, $N_2 \circ \alpha_1 = n_1$, ce qui entraîne que $N_2$
est $n_1 \otimes \mathrm{id}_{A_2}$ sur l'image de $A_1 \otimes A_2$ (vérifié sur les
$x_1 \otimes x_2$, où les deux membres valent $n_1(x_1)x_2^{2}$, puis sur les formes
polaires) ; et (2) entraîne (1)\footnote{Page 111, la flèche du diagramme est lue
$T_1$, et page 112 le premier facteur se lit $t_1(x_1)$, le changement de base
« $k \to A_1$ » ; on lit $T_2$, $n_1(x_1)$ et $k \to A_2$, comme l'exige le texte.} ;
\item[(3)] $\sigma_1\sigma_2\sigma_3 = 1$, qui force les $\sigma_i$ à commuter (deux
involutions commutent si et seulement si leur produit est une involution) et à
former un groupe de Klein. Explicitement, $2x = T_1(x) + T_2(x) + T_3(x) -
t_1T_1(x)$, et, par symétrie, $t_1T_1 = t_2T_2 = t_3T_3 =: T_{A/k}$ quand $k \to A$
est injectif, d'où
\[
  2x = T_1(x) + T_2(x) + T_3(x) - T_{A/k}(x).
\]
\end{enumerate}
Les relations entre normes en découlent : avec $x_0 = x$, $x_j = \sigma_j(x)$,
$n_i(N_i(x)) = x_0x_1x_2x_3 = N_{A/k}(x)$ pour chaque $i$, et
$N_1(x)N_2(x)N_3(x) = x^{2}N_{A/k}(x)$\footnote{Page 114, on lit $x_j = \sigma_j(x)$
où la page écrit $T_j(x)$. Le haut de la page, barré, énonce ces formules comme
attentes.}.

\emph{L'action du groupe pair.} Sur $A_1 \otimes A_2 \otimes A_3$ opère $\mathbf{F}_2^{3}$,
mais il n'y a pas en général d'action compatible sur $A$ : un $\tau_1$ qui fixe $A_2$
et $A_3$ fixe $A$ tout entier si $A_2$, $A_3$ l'engendrent. Le sous-groupe
$\Gamma = \operatorname{Ker}(\mathbf{F}_2^{3} \xrightarrow{\Sigma} \mathbf{F}_2)$, lui, opère, $\tau_i$
agissant sur $A$ par $\sigma_i$. Si les $A_i$ sont les algèbres de revêtements
$X_i$, de produit $B = A_1 * A_2 * A_3$, on a $B \to (A_1 \otimes A_2 \otimes A_3)^{\Gamma}$,
isomorphisme hors des points de caractéristique $2$ où deux au moins des $A_i$ ne
sont pas étales. La composée $f : B \to A_1 \otimes A_2 \otimes A_3 \to A$ tombe dans
$A^{\Gamma}$ ; si $A^{\Gamma} = k\cdot 1_A$, elle définit une \emph{section} $\varepsilon$
de $X_1 * X_2 * X_3$. Et si $2$ est inversible, $A^{\Gamma} = k$ : $\sigma_i(x) = x$
pour tout $i$ donne $6x = 2x + T_{A/k}(x)$, donc $x = \tfrac{1}{4}T_{A/k}(x)$. C'est la
version intrinsèque du fait que, dans $k(\sqrt{a}, \sqrt{b})$, les trois
sous-extensions $k(\sqrt{a})$, $k(\sqrt{b})$, $k(\sqrt{ab})$ ont un produit trivial.

\textbf{Proposition} (page 118). Si les $A_i$ sont étales, $k \to A$ injectif,
$A_1 \otimes A_2 \otimes A_3 \to A$ surjectif et $A^{\Gamma} = k$, alors
$\mathcal{X} = \operatorname{Spec} A \to X_i \times X_j$ est un isomorphisme pour chaque
couple, $\mathcal{X}$ est fini étale sur $X$, et la situation se déduit par torsion
de la situation standard : trois ensembles à deux éléments $I_1, I_2, I_3$, une
trivialisation $\varepsilon$ de $I_1 * I_2 * I_3$, et $J = \{(a_i) : a_1 * a_2 * a_3 =
\varepsilon\}$. Réciproquement une telle situation tordue est biquadratique étale,
avec $A^{\Gamma} = k$.

\textbf{Lemme} (pages 119 et 120). Soient $\Gamma$ opérant transitivement sur un ensemble
$J$, et $\mathcal{X} \subset X_J$ un sous-schéma fermé stable par $\Gamma$ tel que
$k \to \Gamma(\mathcal{X}, \mathcal{O}_{\mathcal{X}})$ soit injectif. Alors $\mathcal{X} =
X_J$. En effet $\mathcal{X} \cap X_j$ est défini par un idéal $\mathfrak{J}_j$ de $k$,
indépendant de $j$ par transitivité, et $\bigcap_j \mathfrak{J}_j$ est le noyau de
$k \to \Gamma(\mathcal{X}, \mathcal{O})$, nul\footnote{La page écrit la condition b)
avec l'anneau de $X_J$, $k^{J}$ ; l'application $k \to k^{J}$ est toujours injective,
et la démonstration utilise l'anneau de $\mathcal{X}$.}.

« Ce résultat vaut sans restrictions de caractéristiques, mais il concerne la
situation standard étale, \emph{très} restrictive. » Il veut montrer que pour trois
revêtements quadratiques quelconques et une section $\varepsilon$ de $X_1 * X_2 * X_3$,
$\mathcal{X} = \pi^{-1}(\varepsilon) \subset X_1 \times_X X_2 \times_X X_3$ porte une structure
biquadratique au-dessus des $X_i$, ce qui généraliserait celle de $X_1 \times X_2$
au-dessus de $X_1$, $X_2$ et $X_3 = X_1 * X_2$. La page 121 en esquisse la variante
universelle --- au-dessus de $\mathfrak{X}_0 = X_1 * X_2 * X_3$, l'algèbre
$A_1 \otimes A_2 \otimes A_3$ est de rang $4$ sur $B_0 = A_1 * A_2 * A_3$ --- et bute :
la multiplication n'est ni finie localement libre ni plate. Le dossier n'y revient
pas\footnote{Les lettres de la page 121 sont incohérentes ($a_1, a_2, a_3 \mapsto x_1
\wedge \alpha_2 \wedge x_3$), et l'indice du produit fibré se lit $\mathfrak{X}$ ou
$\mathfrak{X}_0$.}.

\pagerange{123}{191}
\section*{III. Extensions $\chi$-scindées quelconques (pages 123 à 191)}

La troisième partie reprend le formalisme pour des extensions
$0 \to L \xrightarrow{\alpha} E \xrightarrow{\psi} M \to 0$ dont le quotient $M$ n'est
plus $\mathcal{O}$. Ce sont des calculs rapides, sur les rectos de feuillets dont les
versos sont dactylographiés ; la prose de liaison est rare et souvent illisible,
les formules portent l'argument.

\pagerange{123}{131}
\subsection*{$\chi$-scindages et coordonnées (pages 123 à 131)}

Un \emph{$\chi$-scindage} d'une telle extension est un couple $(\pi, \varpi)$,
$\pi : E \to L$, $\varpi : M \to E$, avec
\[
  \pi\alpha = \chi\,\mathrm{id}_L, \qquad \psi\varpi = \chi\,\mathrm{id}_M, \qquad
  \alpha\pi + \varpi\psi = \chi\,\mathrm{id}_E
\]
(page 169) ; pour $M = \mathcal{O}$, c'est la $\chi$-trivialisation de la deuxième
partie. Un scindage ordinaire $p : E \to L$ donne $E \simeq M \oplus L$, et
$b = \pi - \chi p$, nul sur $L$, est un homomorphisme $b : M \to L$. Pour deux
extensions, le produit est $\mathcal{E}_p = (M_1 \otimes M_2) \oplus (L_1 \otimes L_2)$ en
coordonnées, avec la $\chi$-rétraction $(\eta, \xi) \mapsto \chi\xi + (b_1 \otimes b_2)(\eta)$.
Changer de scindage, $p'_i = p_i - u_i$ avec $u_i : M_i \to L_i$, change $b_i$ en
$b_i + \chi u_i$, et l'identification $\mathcal{E}_p \to \mathcal{E}_{p'}$ est la
transvection $(\eta, \xi) \mapsto (\eta, \xi - u\eta)$ avec
\[
  u = b_1 \otimes u_2 + u_1 \otimes b_2 + \chi\, u_1 \otimes u_2 ;
\]
on a bien $b'_1 \otimes b'_2 = b_1 \otimes b_2 + \chi u$. Pour que ces identifications se
recollent, il faut la relation de cocycle $u'' = u + u'$ pour deux changements
successifs, $u'$ étant calculé avec les $b'_i$. La page 131 l'entreprend et
s'arrête ; elle est vraie :
\[
  u'' = u + (b_1 + \chi u_1) \otimes u'_2 + u'_1 \otimes (b_2 + \chi u_2) + \chi\, u'_1 \otimes u'_2 = u + u'.
\]\footnote{La
  vérification est de l'édition. Les signes de $b'_i = b_i \pm \chi u_i$ varient d'une
  ligne à l'autre des pages 129 à 133, et la page 131 écrit $b''_2 = b_2 + \chi u'_2$ ;
  on fixe $b'_i = b_i + \chi u_i$, $b''_i = b'_i + \chi u'_i$, avec lesquels la relation
  est exacte.}
Le produit $E_1 * E_2$ existe donc par recollement.

\emph{Le produit tiré} (pages 125 et 127). Une seconde construction, duale, donne
en coordonnées la même extension, mais avec le $\chi$-scindage défini par $-b$ au
lieu de $b$ ; un isomorphisme de l'une sur l'autre est une transvection $T_v$ avec
$b + \chi v = -b$, c'est-à-dire $v = -\tfrac{2}{\chi}b$. Il existe donc, canonique,
dès qu'il y a $\gamma \in k$ avec $\gamma\chi = 2$ : $v = -\gamma b$, et le carré de la page
125 commute (« commute ! »). Pour $\chi = 2$, $\gamma = 1$\footnote{La page 125 note
$\vee$ et $\wedge$ les deux produits ; on les note $E_1 * E_2$ et $E_1 *^{\varpi} E_2$. La
page 127 note $\sigma$ l'élément qu'on appelle $\gamma$, comme le font les pages 141 et
suivantes.}.

\pagerange{133}{139}
\subsection*{$\lambda$-rétractions et produit des morphismes (pages 133 à 139)}

Les signes « étranges » qui apparaissent page 133 ($u_1 * u_2 = -\chi\, u_1 \otimes u_2$
quand on passe aux morphismes) « sont dus au signe $-\mathrm{id}$ » de l'autodualité.
Notons $\underline{\mathrm{Hom}}^{!}(E, L)$ le Module des $f : E \to L$ dont la
restriction à $L$ est une homothétie $\lambda\,\mathrm{id}_L$ --- des
« $\lambda$-rétractions » ; c'est une extension de $\mathcal{O}$ par $L$, et pour $M =
\mathcal{O}$,
\[
  \Phi : E \xrightarrow{\ \sim\ } \underline{\mathrm{Hom}}^{!}(E, L), \qquad
  \Phi(x)(y) = \psi(x)\,y - \psi(y)\,x ,
\]
est l'isomorphisme $g$ de la page 95, qui induit $-\mathrm{id}_L$ (pages 133 et 135).
Le produit s'étend aux $\lambda$-rétractions : pour $f_i$ une $\lambda_i$-rétraction de
$E_i$, la $\lambda_1\lambda_2$-rétraction $f_1 * f_2$ de $E_1 * E_2$ est donnée par
\[
  \boxed{\ (f_1 * f_2)(x_1 * x_2) = \lambda_1\pi_1(x_1) \otimes f_2(x_2) + f_1(x_1) \otimes \lambda_2\pi_2(x_2)
  - \chi\, f_1(x_1) \otimes f_2(x_2)\ }
\]
et $(f_1 * f_2)|_{L_1 \otimes L_2} = \lambda_1\lambda_2\,\mathrm{id}$. Pour $f_i = \pi_i$
($\lambda_i = \chi$), on retrouve $\pi_1 * \pi_2 = \chi\,\pi$ ; sur $x_i = u_i \in L_i$, le
second membre vaut $\chi\lambda_1\lambda_2\, u_1 \otimes u_2$, ce qui est bien la valeur de
$f_1 * f_2$ sur $u_1 * u_2 = \chi\, u_1 \otimes u_2$\footnote{La page 135 écrit
$f_1(x_2)$ dans le deuxième terme ; la page 137 corrige en $f_1(x_1)$. Le dernier
calcul de la page 139, où apparaît « $2 - \chi = 1$ », est surchargé à l'endroit
décisif ; on ne le reproduit pas. La vérification sur $L_1 \otimes L_2$ est de
l'édition.}. La page 139 retrouve aussi la rétraction $p$ de $E_1 * E_2$ déduite de
scindages $p_1, p_2$ : $p(x_1 * x_2) = x_1 * p_2(x_2) + p_1(x_1) * x_2 - p_1(x_1) * p_2(x_2)$.

\pagerange{141}{157}
\subsection*{L'opérateur $F_{11}$ et la condition $\gamma\chi = 2$ (pages 141 à 157)}

Il cherche à identifier les deux produits par un endomorphisme $F$ de
$E_1 \otimes E_2$, dont il connaît deux restrictions : sur $E_1 \otimes_0 E_2$, $F_{11}$ doit
se factoriser par $\pi : E_1 \otimes_0 E_2 \to L_1 \otimes L_2$ ; et sa composée avec la
projection sur $E_1 \otimes^{\circ} E_2$ doit se factoriser par $\psi_1 \otimes \psi_2$. Il
pose $\pi'_1 = \pi_1 \otimes \mathrm{id}$, $\pi'_2 = \mathrm{id} \otimes \pi_2$ (relevés dans
$E_1 \otimes E_2$), cherche $F_{11}$ de la forme
\[
  F_{11} = \alpha + \beta(\pi'_1 + \pi'_2) + \gamma\,\pi'_1 \pi'_2,
\]
et trouve, par la table de $\pi'_i$ sur les quatre morceaux $L_1 \otimes L_2$,
$L_1 \otimes M_2$, $M_1 \otimes L_2$, $M_1 \otimes M_2$ (pages 153 et 163), les conditions
\[
  \beta = 1 - \gamma\chi, \qquad \alpha = -\chi(1 - \gamma\chi)
\]
(pages 157 et 163, où la constante s'appelle tour à tour $\gamma$ et $\delta$). Les deux
conditions connues sont compatibles et laissent une indétermination, un
homomorphisme $M_1 \otimes M_2 \to L_1 \otimes L_2$ : « C'est normal […] on a
l'indétermination des choix ». La normalisation qui fait de
$\sigma_i = \mathrm{id} - \gamma\pi_i$ une involution est $\gamma\chi = 2$ : comme
$\pi_i^{2} = \chi\pi_i$ (page 149), $(\mathrm{id} - \gamma\pi_i)^{2} = \mathrm{id} + \gamma(\gamma\chi -
2)\pi_i$. Alors $\beta = -1$, $\alpha = \chi$, et pour $\chi = 2$, $\gamma = 1$ :
\[
  F_{11}(x_1 \otimes x_2) = 2\,x_1 \otimes x_2 - \pi_1(x_1) \otimes x_2 - x_1 \otimes \pi_2(x_2) + \pi_1(x_1) \otimes \pi_2(x_2)
  = x_1 \otimes x_2 + \bar{x}_1 \otimes \bar{x}_2 ,
\]
avec $\bar{x}_i = x_i - \pi_i(x_i)$ : c'est l'opérateur $\xi \mapsto \xi + \bar{\xi}$ des pages
98 à 104, « o.k. » (page 155)\footnote{Pages 147 à 151, il essaie aussi
$\sigma_i = \lambda + \mu\pi_i$, d'où $\lambda^{2} + 1 = \chi$, $\lambda\mu = -1$, $\mu^{2} =
\gamma$ ; l'essai n'est pas poursuivi. La page 143 esquisse, pour des
endomorphismes $f_i$ de $E_i$ qui stabilisent $L_i$, l'endomorphisme induit sur le
produit ; la page 145 reprend en matrices les scindages de la page 129. La page 151
écrit « ok si $\chi = \pm 2$ » et s'interrompt.}.

\pagerange{159}{169}
\subsection*{Produits poussé et tiré ; dualité (pages 159 à 169)}

La page 169 met en place les deux constructions. Le \emph{produit poussé}
$E_1 * E_2$ est la somme amalgamée de $E_1 \otimes E_2$ et de $L_1 \otimes L_2$ le long de
$(\pi_1 \otimes \mathrm{id}, \mathrm{id} \otimes \pi_2) : E_1 \otimes_0 E_2 \to L_1 \otimes L_2$, où
$E_1 \otimes_0 E_2$ est la somme amalgamée de $E_1 \otimes L_2$ et $L_1 \otimes E_2$ le long de
$L_1 \otimes L_2$. Le \emph{produit tiré} $E_1 *^{\varpi} E_2$ est le produit fibré de
$E_1 \otimes E_2$ et de $M_1 \otimes M_2$ au-dessus de $E_1 \otimes^{\circ} E_2 = (E_1 \otimes M_2)
\times_{M_1 \otimes M_2} (M_1 \otimes E_2)$, le long de $(\varpi_1 \otimes \mathrm{id}, \mathrm{id} \otimes
\varpi_2)$. Ce sont deux extensions de $M_1 \otimes M_2$ par $L_1 \otimes L_2$, munies de
$\chi$-scindages, et il y a un isomorphisme canonique de l'une sur l'autre, sous la
condition $\gamma\chi = 2$ de la page 127 :

\begin{tikzcd}[column sep=small, nodes={font=\small}]
0 \arrow[r] & E_1 \otimes_0 E_2 \arrow[r] \arrow[d] & E_1 \otimes E_2 \arrow[r] \arrow[d] & M_1 \otimes M_2 \arrow[r] \arrow[d, no head, "=" description] & 0 \\
0 \arrow[r] & L_1 \otimes L_2 \arrow[r] \arrow[d, no head, "=" description] & E_1 * E_2 \arrow[r] \arrow[d, "\wr"] & M_1 \otimes M_2 \arrow[r] \arrow[d, no head, "=" description] & 0 \\
0 \arrow[r] & L_1 \otimes L_2 \arrow[r] \arrow[d, no head, "=" description] & E_1 *^{\varpi} E_2 \arrow[r] \arrow[d] & M_1 \otimes M_2 \arrow[r] \arrow[d] & 0 \\
0 \arrow[r] & L_1 \otimes L_2 \arrow[r] & E_1 \otimes E_2 \arrow[r] & E_1 \otimes^{\circ} E_2 \arrow[r] & 0
\end{tikzcd}

La composée verticale du milieu est l'opérateur $F_{11}$\footnote{Le diagramme de la
page 167 porte « ? » sur la flèche du milieu, et un trait barre la page entre ses
deux premières lignes. On le lit avec la flèche que les pages 141 à 157
construisent. La page 169 dit l'isomorphisme « de façon évidente » ; l'hypothèse
$\gamma\chi = 2$ est de l'édition, d'après la page 127.}. Les deux constructions se
correspondent par transposition (page 159) : le dual d'une extension
$0 \to L \to E \to M \to 0$ est $0 \to M^{\vee} \to E^{\vee} \to L^{\vee} \to 0$, la
transposée d'une $\chi$-rétraction est une $\chi$-section, et, pour $M_i$ localement
libres de rang fini, $E_1^{\vee} *^{\varpi} E_2^{\vee} \simeq (E_1 * E_2)^{\vee}$. Pour
$M = \mathcal{O}$ et $\chi = 2$, c'est la dualité entre l'algèbre $A_1 * A_2$, poussée
le long de la trace (pages 99 et 100), et l'extension $E_1 * E_2$\footnote{La page 159
note $L'$ le quotient, et un exposant de $L$, en forme de « $4$ », est transcrit
« ! » ; la page 161 pose $\pi' = \pi - (\chi - 1)\,\mathrm{id}$, nul ou presque sur
$L$, dans un calcul qui se termine par un point d'interrogation.}.

\pagerange{171}{191}
\subsection*{Retour aux algèbres quadratiques : l'accouplement (pages 171 à 191)}

Pour des algèbres quadratiques $A_i$ et leurs duaux $E_i$, il récrit les deux
constructions côte à côte :
\[
  A = (A_1 \otimes A_2) \amalg_{A_1 \otimes_0 A_2} \mathcal{O}, \qquad
  E = (E_1 \otimes E_2) \amalg_{E_1 \otimes_0 E_2} (L_1 \otimes L_2),
\]
la première le long de $(\mathrm{Tr}_1, \mathrm{Tr}_2)$, la seconde le long de
$(\pi_1 \otimes \mathrm{id}, \mathrm{id} \otimes \pi_2)$, avec les flèches $A_1 \otimes A_2 \to A
\to A_1 \otimes A_2$ et $E_1 \otimes E_2 \to E \to E_1 \otimes E_2$ dont les composées sont
$f \mapsto f + \bar{f}$, $x \mapsto x + \bar{x}$, transposées l'une de l'autre :
${}^{t}(\mathrm{id} + \sigma_1 \otimes \sigma_2) = \mathrm{id} + {}^{t}\sigma_1 \otimes {}^{t}\sigma_2$.
L'application $A \to A_1 \otimes A_2$ est injective sauf en caractéristique $2$ là où
$b_1 = b_2 = 0$ (page 171).

\emph{L'accouplement} $A \times E \to \mathcal{O}$ sur les deux sommes amalgamées est
donné (pages 179 et 183) par
\[
  \langle\lambda, u\rangle = 0, \qquad \langle f, u\rangle = \varphi(f)\cdot u, \qquad
  \langle\lambda, x\rangle = \lambda\,\psi(x),
\]
pour $\lambda \in \mathcal{O}$, $u \in L_1 \otimes L_2$, $f \in A_1 \otimes A_2$ d'image $\varphi(f)$
dans $L_1^{\vee} \otimes L_2^{\vee}$, $x \in E_1 \otimes E_2$, et par l'accouplement produit
sur $(A_1 \otimes A_2) \times (E_1 \otimes E_2)$ ; il faut quatre compatibilités pour qu'il
passe aux sommes amalgamées. La page 177 en écrit trois et conclut « OK » ; la
quatrième n'est pas écrite\footnote{La page 179 écrit $\langle\lambda, x\rangle =
\lambda\,\pi(x)$ ; la page 183 écrit $\lambda\,\psi(x_1 \otimes x_2)$, qui est le bon. Les
pages 181 à 187 essaient plusieurs écritures de $A$ et de $A^{\vee}$ comme
quotients, en partie barrées.}.

\emph{Le générateur} (pages 175 et 191). Avec les bases duales $(1, u)$ de $A$ et
$(x, e)$ de $E$, $x = x_1 * x_2$, $e = e_1 \otimes e_2$, $\langle 1, x\rangle = \langle u,
e\rangle = 1$, $\langle 1, e\rangle = \langle u, x\rangle = 0$, les coordonnées de
$u_1 * u_2$ sont $\langle u_1 * u_2, x\rangle = b_1b_2$ et $\langle u_1 * u_2, e\rangle = 1$,
d'où
\[
  \boxed{\ u_1 * u_2 = u + b_1b_2\ }, \qquad i(u) = 2u_1u_2 + b_2u_1 + b_1u_2 ,
\]
ce qui est la formule de la page 102, retrouvée par la dualité. Les pages 189 à 191
recommencent le calcul en coordonnées --- $e_i$ base de $L_i$, $b_i = \beta_ie_i$,
$c_i = \gamma_ie_i^{\otimes 2}$, $u_i^{2} + \beta_iu_i + \gamma_i = 0$ ---, notent $\mathrm{Tr}(\lambda u
+ \mu) = -\lambda\beta + 2\mu$ et $N(\lambda u + \mu) = \lambda^{2}c - \lambda\mu\beta + \mu^{2}$, et
posent l'équation $u^{2} + \beta u + \gamma = 0$ du générateur du produit. Le dossier
s'arrête là, sans conclusion visible\footnote{Page 189, $\lambda$ et $\mu$ sont les
coefficients de $u$ et de $1$, l'inverse de la page 175 ; et $u$ est défini par
$\langle u, e_1 \otimes e_1\rangle = 1$, pour $e_1 \otimes e_2$. Page 191, « $u = 2uu' +$ »
s'interrompt après le signe $+$.}.

\end{document}
